\documentclass[10pt,a4paper]{amsart}
\usepackage[margin=19mm]{geometry}
\usepackage{amsmath,amssymb,mathtools}
\usepackage[scr=rsfs,cal=euler]{mathalfa}
\usepackage{xfrac}
\usepackage[unicode,hidelinks,bookmarks=false]{hyperref}
\newcommand{\ie}{i.e.\ }
\newcommand{\cf}{cf.\ }
\newcommand{\resp}{respectively\ }
\newcommand{\Subsection}{\subsection}
\providecommand{\nobreakdash}{\nobreak\hbox{-}}
\setlength{\emergencystretch}{2em}
\multlinegap0pt
\usepackage{mathrsfs,amsopn}
\usepackage{xy}
\usepackage{dsfont}
\usepackage{color}
\usepackage{esint}
\usepackage{mathtools}
\mathtoolsset{showonlyrefs}
\usepackage{slashed}
\usepackage{tikz-cd}
\usepackage{faktor}
\newcommand{\tnorm}[1]{%
  \left\vert\kern-0.3ex\left\vert\kern-0.3ex\left\vert #1 \right\vert\kern-0.3ex\right\vert\kern-0.3ex\right\vert
}
\usepackage{todonotes}
\newcommand{\todoRuijun}[2][]{\todo[author=Ruijun, color=red!10, #1]{#2}} 
\newcommand{\Bd}{\mathcal{B}}
\newcommand{\p}{\partial}
\newcommand{\R}{\mathbb{R}}
\newcommand{\C}{\mathbb{C}}
\newcommand{\T}{\mathbb{T}}
\newcommand{\QH}{\mathbb{H}}
\newcommand{\RS}{\Sigma}
\newcommand{\D}{\slashed{D}}
\newcommand{\Mnabla}{\overline{\nabla}} % Levi-Civita of ambient space
\newcommand{\Penr}{\slashed{P}}
\newcommand{\pd}{\slashed{\partial}}
\newcommand{\dd}{\mathop{}\!\mathrm{d}}
\newcommand{\snabla}{\slashed{\nabla}} % spin connection
\newcommand{\sph}{\mathbb{S}}
\newcommand{\shape}{A}
\newcommand{\eps}{\varepsilon}
\newcommand{\pD}{\mathbf{D}}
\newcommand{\G}{\slashed{G}} % Green operator for Dirac
\newcommand{\Abracket}[1]{\left<#1\right>} % angle bracket 
\newcommand{\parenthesis}[1]{\left(#1\right)} % round bracket
\newcommand{\braces}[1]{\left\{#1\right\}} % curly bracket
\newcommand{\cH}{\mathcal{H}}
\newcommand{\bZ}{\mathbb{Z}}
\newcommand{\cW}{\mathcal{W}}
\newcommand{\cO}{\mathcal{O}}
\DeclareMathOperator{\Aut}{Aut}
\DeclareMathOperator{\Capa}{cap} % capacity
\DeclareMathOperator{\Cl}{Cl}
\DeclareMathOperator{\CCl}{\mathbb{C}l}
\DeclareMathOperator{\codim}{codim}
\DeclareMathOperator{\Coker}{Coker}
\DeclareMathOperator{\diag}{diag}
\DeclareMathOperator{\dist}{dist}
\DeclareMathOperator{\Dom}{Dom}
\DeclareMathOperator{\dv}{\dd{vol}}
\DeclareMathOperator{\ds}{\dd{s}}
\DeclareMathOperator{\dx}{dx}
\DeclareMathOperator{\diverg}{div}
\DeclareMathOperator{\Eigen}{Eigen}
\DeclareMathOperator{\End}{End}
\DeclareMathOperator{\grad}{grad}
\DeclareMathOperator{\GL}{GL}
\DeclareMathOperator{\gAPS}{gAPS}
\DeclareMathOperator{\Hom}{Hom}
\DeclareMathOperator{\Ind}{Index}
\DeclareMathOperator{\id}{Id}
\newcommand{\II}{\mathrm{I\hspace*{-0.4ex}I}} % Second fundmantal form of N
\DeclareMathOperator{\Ker}{Ker}
\DeclareMathOperator{\mass}{mass}
\DeclareMathOperator{\Mat}{Mat} % matrix ring
\DeclareMathOperator{\Mtp}{Multi} % multiplicity
\DeclareMathOperator{\Orth}{O}
\DeclareMathOperator{\Pin}{Pin}
\DeclareMathOperator{\Prin}{Prin} % equiv classes of principal bundles
\DeclareMathOperator{\proj}{Pr}
\DeclareMathOperator{\Ran}{Ran} % range of an operator 
\DeclareMathOperator{\rank}{rank}
\DeclareMathOperator{\Ric}{Ric}
\DeclareMathOperator{\Span}{Span}
\DeclareMathOperator{\Scal}{Scal}
\DeclareMathOperator{\Sk}{Sk} % Skeleton
\DeclareMathOperator{\Spect}{Spect}
\DeclareMathOperator{\St}{St}
\DeclareMathOperator{\Spin}{Spin}
\DeclareMathOperator{\spt}{spt}
\DeclareMathOperator{\SO}{SO}
\DeclareMathOperator{\tr}{tr}
\DeclareMathOperator{\vol}{vol}
\DeclareMathOperator{\TRC}{TRC}
\DeclareMathOperator{\up}{up}
\DeclareMathOperator{\down}{down}
\newtheorem{thm}{Theorem}[section]
\newtheorem{dfn}[thm]{Definition}
\newtheorem{lemma}[thm]{Lemma}
\newtheorem{prop}[thm]{Proposition}
\newtheorem{cor}[thm]{Corollary}
\newtheorem{fact}[]{FACT}
\newtheorem{rmk}[thm]{Remark}
\newtheorem{eg}{Example}
\newtheorem{theorem}{Theorem}
\renewcommand*{\thetheorem}{\Alph{theorem}}
\begin{document}
\title[Stationary periodic solutions for nonlinear Dirac equations ]{Stationary periodic solutions for nonlinear Dirac equations with non-coercive nonlinearity II: splitting}
\author{Ruijun Wu; Fuping Zhang}
\date{}
\maketitle
\noindent\emph{Source and licence.} Stationary periodic solutions for nonlinear Dirac equations with non-coercive nonlinearity II: splitting. Version arXiv:2608.17768v1 (CC BY 4.0). This material is distributed under the Creative Commons Attribution 4.0 International licence, http://creativecommons.org/licenses/by/4.0/, as recorded in the arXiv OAI licence metadata for this exact version and shown on the abstract page https://arxiv.org/abs/2608.17768v1. Full rights notice: licenses/arxiv-2608.17768-CC-BY-4.0.txt.\par\smallskip
\noindent\textit{Editorial scope.} Counted unit: the original \texttt{thm} environment \texttt{thm:uniform estimate} at source lines 1688--1693 together with its complete proof at lines 1695--1959; the delivered document keeps source lines 158--1960 so that every referenced label and every citation resolves inside it. Native editable source is the paper's own concatenated TeX; the AMS wrapper is editorial formatting only. No publisher class, upstream script or unrelated artwork is redistributed.
\par\smallskip\noindent\textit{Reference note.} This article ships no compiled bibliography (biblatex); the entries delivered here are composed from the author's own BibTeX (.bib) fields, with no field invented and no entry dropped.
\begin{align}\label{eq:NDE-spatial}
    -\sum_{j=1}^3 i\gamma^0\gamma^j\p_j\psi + m\gamma^0\psi =a\psi + \p F
\end{align}
on a flat three-dimensional torus~$\T^3$, where~$m>0$ stands for the mass,~$a>0$ is a temporal frequency,~$F\colon \T^3\times \C^4\to \R$ is the nonlinearity standing for the self-coupling interaction and~$\p F(\psi(x))\equiv \p_\psi F(x,\psi(x))$ denotes the Nemytskii operator associated to~$F$.
The~$\gamma^\mu$'s are given as follows.
The matrix~$\gamma^0$ takes the form
\begin{align}
    \gamma^0= \begin{pmatrix} I_2 & 0 \\ 0 & -I_2\end{pmatrix}
    =\begin{pmatrix}
        1 & & &  \\ & 1 & & \\ & & -1 & \\ & &  & -1
    \end{pmatrix}
\end{align}
and induces a decomposition of the spinors into~$\pm1$ eigenspaces of~$\gamma^0$.
As for the others, we denote by
\begin{align}
    {\sigma^1} = \left( {\begin{array}{*{20}{c}}
			{0}&1 \\
			1&{0}
	\end{array}} \right),\;{\sigma^2} = \left( {\begin{array}{*{20}{c}}
			{0}&{ - i} \\
			i&{0}
	\end{array}} \right),\;{\sigma^3} = \left( {\begin{array}{*{20}{c}}
			1&{0} \\
			{0}&{ - 1}
	\end{array}} \right)
\end{align}
for the Pauli matrices, and the gamma matrices are
\begin{align}
    {\gamma^k} = \left( {\begin{array}{*{20}{c}}
			{0}&{{\sigma^k}} \\
			{{-\sigma^k}}&{0}
	\end{array}} \right),\quad\text{for}\quad k= 1, 2, 3.
\end{align}
Note that they satisfy the following Clifford relations: for~$j,k\in\braces{1,2,3}$,
\begin{align}
    \sigma^j\sigma^k+\sigma^k\sigma^j= 2\delta^{jk}I_2, & &
    \gamma^j\gamma^k + \gamma^k\gamma^j= -2\delta^{jk}I_4.
\end{align}
Moreover,
\begin{align}
    \gamma^0\gamma^k + \gamma^k \gamma^0 =0, & &
    (\gamma^0)^2=I_4.
\end{align}
Then~$\sum_{\mu} i\gamma^\mu \p_\mu$ is the four-dimensional space-time Dirac operator.


Equations of such forms have been widely studied in physics, chemistry and mathematics, see~\cite{Esteban2002overview,Ranada1983Classical,Thaller1992Dirac} and the references therein. 
They are governed by a general simple Dirac equation in the Minkowski spacetime 
\begin{align}\label{eq:NDE-spacetime}
    \sum_{\mu=0}^3 i\gamma^\mu\p_\mu \Psi - m \Psi +\gamma^0 \p_\Psi F(\cdot,\Psi)=0, \qquad \mbox{ in } \; \R^{1+3}. 
\end{align}
Here~$x^0$ stands for the time coordinate, and the equation is hyperbolic. 
The equation can be reduced to~\eqref{eq:NDE-spatial} if we consider stationary solutions, namely spinors of the form
\begin{align}\label{eq:ansatz}
    \Psi(x^0, \vec{x}) = e^{-iax^0}\psi(\vec{x})
\end{align}
for some~$a\in \R\setminus\braces{0}$. 
Note that~$|\Psi|^2=|\psi|^2$, so~$\Psi$ and~$\psi$ possess the same probability density, and as time evolves the state is essentially unchanged in shape. 
Additional hypotheses on~$F$ are required before one can really treat~\eqref{eq:NDE-spacetime} as an equation on~$\R^3$, for example~$F$ should not depend on the time coordinate explicitly, and~$F$ is invariant under phase rotation:
\begin{align}
    F(x, e^{i\theta}\psi)=F(x,\psi), \quad \forall e^{i\theta} \in U(1). 
\end{align}
 
Now we restrict to consider~\eqref{eq:NDE-spatial}.
A large part of the variational literature deals with nonlinearities of Ambrosetti--Rabinowitz type, namely, the nonlinearities consist of functions in the Euclidean amplitude~$|\psi|$ and small perturbations. 
In that setting the main difficulties arise from the strong indefiniteness of the functional, failure of Palais--Smale conditions, and the competition with the nonlinear term, see for example~\cite{BartschDing2006Solutions,DingRuf2008Solutions,YangDing2012Stationary}.
On general compact spin manifolds, related variational problems have also been studied, see for instance~\cite{Ammann2003Habil, BartschXu2021spinorial, DingLiXu2016bifurcation, DingXu2025localized, Isobe2011nonlinear, Isobe2020Morse}. 



Another class of nonlinearities consists of those involving Lorentz type scalars such as~$\bar{\psi}\psi\equiv \Abracket{\gamma^0\psi,\psi}$. 
They arise naturally in standard fermionic models.
In~\cite{Soler1970Classical} and the even earlier~\cite{Finkelstein1951Non-linea}, the authors considered the nonlinearity~$ (\bar{\psi}\psi)^2$ in the fermionic spinor functional. 
For this reason we refer to the nonlinearities of the form~$|\bar{\psi}\psi|^\nu$ as Soler-type nonlinearities; sometimes we also allow for a positive smooth function multiplied to this nonlinearity, standing for certain nonconstant field, as well as small perturbations added, see the examples in~\cite{WZ2026Stationary}. 

Writing the spinor as~$\psi=(\psi_{\up},\psi_{\down})$ as in~\eqref{eq:up-down}, we see clearly the degeneracy of this Lorentz scalar:
\begin{align}
    \bar{\psi}\psi=\Abracket{\gamma^0\psi,\psi}
    =|\psi_{\up}|^2 - |\psi_{\down}|^2.
\end{align}
Thus~$\bar{\psi}\psi$ may stay small even if~$|\psi|$ is large. 
Such a degeneracy is inherited by the nonlinearity~$F$: 
in the Lorentz null cone
\begin{align}
    \mathcal{N}=\braces{\psi\in\C^4 \mid \bar{\psi}\psi=0}, 
\end{align}
the nonlinearity~$F$ may be small, but this cone contains spinors of arbitrary large norm. 
This causes severe difficulties for the existence theory of~\eqref{eq:NDE-spatial}.
The autonomous case in~$\R^3$ has been studied extensively.
Via a radial ansatz, one can reduce to a system of ODEs in two functions, which can be solved by the shooting method as well as techniques from dynamical systems~\cite{BalabaneCazenave1988Existence,Cazenave1986Existence, EstebanSere1995stationary,Merle1988Existence}.


For symmetry reasons it is also natural to consider periodic solutions. 
In~\cite{DingLiu2014periodic} the authors dealt with nonlinearities of Ambrosetti--Rabinowitz type. 
Remark that there are also works studying~\eqref{eq:NDE-spatial} in~$\R^3$ which is non-autonomous but with~$F$ being periodic in the~$x$ variables~\cite{BartschDing2006Solutions,DingLiu2015Periodic,DingLiu2017Periodic}. 
For nonlinearities of Soler type, we recently obtained existence results for~$F$ satisfying additional structural constraints in~\cite{WZ2026Stationary}. 
The purpose of the present paper is to enlarge the frequency regime which admits nonzero solutions.

The main idea is to exploit the product decomposition~$\T^3=\T^2\times\sph^1$.
For a fixed circle mode~$\zeta$, the splitting ansatz
\begin{align}
    \psi=\varphi\otimes e^{i\zeta x^3}
\end{align}
reduces the equation to a nonlinear Dirac equation on~$\T^2$ with the lowest positive eigenvalue of the reduced Dirac operator being~$\sqrt{\zeta^2+m^2}$.
Although the corresponding eigenspace is no longer contained in the~$(+1)$-eigenspace of~$\gamma^0$, it still intersects the Lorentz null cone only at the origin. 
This leads to a quantitative null-cone separation principle. 
Together with a resolvent estimate on the orthogonal complement, this rules out blow-up of the perturbed minimax solutions whose levels tend to zero and whose normalized profiles approach the null cone. 
On the variational side, the perturbed functionals admit linking levels with upper bound~$C^*(a)$ such that 
\begin{align}
    C^*(a)\to 0 \qquad\mbox{as }a\to \sqrt{\zeta^2+m^2}. 
\end{align}
These two ingredients allow us to obtain estimates uniform in both the
coercive perturbation and~$a$ for~$a$ sufficiently close to~$\sqrt{\zeta^2+m^2}$
from below. 
Then we can pass to the zero-perturbation limit and obtain a nonzero solution. 

\ 

We need structure conditions on~$F\colon \T^3\times \C^4\to \R_+$ similar to those in~\cite{WZ2026Stationary} and in~\cite{EstebanSere1995stationary}, which are listed as follows. 
First of all, as already commented, we assume that~$F$ is invariant under phase rotation and does not depend on the third coordinate~$x^3$ in~$\T^3$.
Hence we can view~$F$ as defined on~$\T^2\times\C^4$. 
In addition, the following structural assumptions are imposed: there exist constants~$A_j>0$,~$1\leq j\leq 5$, and~$\nu>1$,~$2<\alpha_1\leq\alpha_2<4$,~$\alpha\in(0,1)$,~$\beta>2$ such that 
\begin{itemize}
    \item[(F1)] $ 0 \leq F(x,\varphi) \leq A_1 \big( |\varphi|^{\alpha_1} + |\varphi|^{\alpha_2} \big)$, $\forall x\in\T^2$ and $\forall \varphi \in \mathbb{C}^4 $;

    \item[(F2)] \( F \in C^{1,\alpha}_{loc} \), $F(x,0) =0 $, $\p_\psi F(x,0)=0$, and
          \( |\p_{\varphi} F(x,\varphi)| \leq A_2 |\varphi|^{\alpha_2 - 1} \) for \( |\varphi| \) large;

    \item[(F3)] \(  \p_\varphi F(x,\varphi)[\varphi]\geq \alpha_1 F(x,\varphi) \); 


    \item[(F4)] \( F(x,\varphi) \geq A_3 |\bar{\varphi}\varphi|^\nu - A_4 \);

    \item[(F5)] for any~$\delta > 0$, there is a~$C_\delta > 0$, such that $\forall \varphi \in \mathbb{C}^4$,~$\forall x\in\mathbb{T}^2$,
        \begin{align}
        |\p_\varphi F(x,\varphi)| 
        \leq A_5 \left( \delta + C_\delta F(x,\varphi)^{\frac{1}{\beta}} \right) |\varphi|.
        \end{align}
\end{itemize}

\begin{rmk}
    Some more comments are in order. 
    \begin{itemize}
        \item[(i)]  By (F1) and (F4) we necessarily have~$\nu\leq \frac{\alpha_2}{2}<2$.

        \item[(ii)] The growth condition in (F2) holds for~$|\varphi|$ large. 
                    But together with (F1) we see that
                    \begin{align}\label{eq:F2-all}
                        |\p F(x,\varphi)|\leq C(1+|\varphi|^{\alpha_2-1}), \qquad \forall (x,\varphi)\in\T^2\times\C^4,  
                    \end{align}
                    for some constant~$C$. 
                    For later convenience we will assume~$A_2$ is large enough that~\eqref{eq:F2-all} holds with~$A_2$. 

        \item[(iii)] Examples of such nonlinearities satisfying (F1-5) include 
                    \begin{align}
                        F(x,\varphi)=h_1(x)|\bar{\varphi}\varphi|^{\frac{\alpha_1}{2}} + h_2(x)|\bar{\varphi}\varphi|^{\frac{\alpha_2}{2}}
                    \end{align}
                    where~$h_1,h_2\colon \T^2 \to \R_+$ are smooth uniformly bounded positive functions, and~$2<\alpha_1 \leq \alpha_2 <4$. 
                    Then we can take a suitable~$\beta\in (2,\frac{\alpha_2}{\alpha_2-2})$. 
                    For more examples we refer to~\cite{WZ2026Stationary}. 
    \end{itemize}
\end{rmk}

We take a lattice~$\Gamma(\ell_1,\ell_2,\ell_3)$ of rank three, and consider the tori 
\begin{align}
    \T^3 =\faktor{\R^3}{\Gamma(\ell_1,\ell_2,\ell_3)}, &  & 
    \T^2 =\faktor{\R^2}{\Gamma(\ell_1,\ell_2)}. 
\end{align}
We refer to Section~\ref{sect:splitting} for explanation of notations. 
\begin{thm}\label{thm:existence-splitting}
    Let~$\zeta\in\bZ\frac{2\pi}{\ell_3}$, ~$\theta=1-\frac{m^2}{36(\zeta^2+m^2)}$, and assume~$F\colon \T^2\times \C^4\to\R_+$ is phase-rotation invariant and satisfies (F1-5). 
    
    There exists~$m_*\in (\theta \sqrt{\zeta^2+m^2}, \sqrt{\zeta^2+m^2})$ such that for each~$a\in(m_*,\sqrt{\zeta^2+m^2})$, the equation~\eqref{eq:NDE-spatial} admits a nonzero~$C^1$ solution~$\psi$ on~$\T^3=\faktor{\R^3}{\Gamma(\ell_1,\ell_2,\ell_3)}$ of the form~$\psi=\varphi\otimes e^{i\zeta x^3}$.
\end{thm}
Note that for~$\zeta\neq 0$, the above solution is automatically nonconstant, at least in~$x^3$. 
Meanwhile, since~$\sqrt{\zeta^2+m^2}>m$, we get existence with~$a$ above the mass threshold~$m$, and as~$\zeta$ varies in~$\bZ\frac{2\pi}{\ell_3}$, we get larger frequency range than that in~\cite{WZ2026Stationary}. 

Our strategy is inspired by the radial ansatz reduction and more closely by~\cite{SireXu2023variational} where the authors directly use the circle factor to reduce the spinorial Yamabe equation to a codimension-one submanifold, turning the critical nonlinearity to a subcritical one. 
For technical reasons we can only deal with subquartic nonlinearities, but this covers the sub-cubic case in~\cite{WZ2026Stationary}. 

\ 

The paper is organized as follows. 
In Section~\ref{sect:splitting} we split out a circle factor and relate the spinors to those on the two-dimensional torus~$\T^2$.
For each~$\zeta$ in the spectrum of the circle that is split out, we reduce~\eqref{eq:NDE-spatial} to~\eqref{eq:NDE-T2} on~$\T^2$. 
The two-dimensional domain allows for nonlinearities of subquartic growth. 
In preparation for later variational analysis, we clarify the eigenvalues and eigenspaces of the Dirac type operators on~$\T^2$ in Section~\ref{sect:spectral}, and make two quantitative estimates concerning spinors away from the lowest positive eigenspace in Section~\ref{sect:key lemma}. 
These generalize the corresponding conclusions in~\cite{WZ2026Stationary}. 
We realize solutions of the reduced equation~\eqref{eq:NDE-T2} as critical points of a functional~$J_\zeta\colon H^{\frac{1}{2}}(\T^2,\C^4)\to\R$ and use a perturbation strategy which dates back to~\cite{EstebanSere1995stationary} and is well developed in~\cite{WZ2026Stationary}.
When the frequency~$a$ is in a left neighborhood of the minimal positive eigenvalue~$\sqrt{m^2+\zeta^2}$ of the reduced Dirac operator, the perturbed functional admits a local level-linking uniformly in~$\eps$.
Using the admissible deformation theory developed in~\cite{WZ2026Stationary}, the local linking geometry leads to a positive minimax level~$\Lambda_1(\zeta,\eps)$ for each~$\eps\in (0,1]$, which can be verified to be a critical level of~$J_{\zeta,\eps}$.
Moreover, we can estimate these critical levels in a fine way. 
At the end of Section~\ref{sect:perturbed solutions} we get a uniform~$H^1$ estimate of such perturbed minimax solutions under the condition that~$a$ is in a sufficiently small left neighborhood of~$\sqrt{m^2+\zeta^2}$. 
Passing to a convergent subsequence, we show that the perturbations in both the equation and the functional vanish in the limit~$\eps_n\to 0^+$. 
The limit is the desired nonzero solution in the main Theorem~\ref{thm:existence-splitting}. 
The regularity can also be verified since the zero order terms are H\"older continuous and the Dirac operator is elliptic of first order. 






\section{Splitting of the torus and reduction of the equation}\label{sect:splitting}

In this section we clarify the relations between the spinors on~$\T^3$ and~$\T^2=\faktor{\T^3}{\sph^1}$, as well as the Dirac operators on these manifolds. 

Let~$\Gamma(\ell_1,\ell_2,\ell_3)=\ell_1\mathbb{Z} \times \ell_2\mathbb{Z} \times \ell_3\mathbb{Z}$ be a product lattice of rank three in~$\R^3$, whose dual lattice is 
\begin{align}
    \Gamma(\ell_1,\ell_2,\ell_3)^*=\frac{1}{\ell_1}\bZ\times \frac{1}{\ell_2}\bZ \times \frac{1}{\ell_3}\bZ.
\end{align}
The periodic functions and spinors can be viewed as defined on the torus
\begin{align}
    \T^3\coloneqq \faktor{\R^3}{\Gamma(\ell_1,\ell_2,\ell_3)} 
    = \sph^1\parenthesis{\frac{\ell_1}{2\pi}} 
    \times \sph^1\parenthesis{\frac{\ell_2}{2\pi}}
    \times \sph^1\parenthesis{\frac{\ell_3}{2\pi}}.
\end{align}
Since this torus admits circle actions (in more than one way), we can split out a circle factor and get 
\begin{align}\label{eq:product-str-mfld}
    \T^3= \T^2\times \sph^1(\frac{\ell_3}{2\pi})
\end{align}
where~$\T^2=\faktor{\R^2}{\Gamma(\ell_1,\ell_2)}$ and~$\Gamma(\ell_1,\ell_2)=\ell_1\bZ\times \ell_2\bZ\subset \R^2$ is a rank two lattice in~$\R^2$.
With a little abuse of notation we will write~$\T^3=\T^2\times\sph^1$. 

In this work we work with spinors which are~$\C^4$ valued functions on~$\T^3$.
This amounts to considering the trivial spin structure on~$\T^3$ and a double of the trivial spinor bundle~$\Sigma \T^3 \oplus\Sigma \T^3=(\T^3\times \C^2)\oplus (\T^3\times\C^2)=\T^3\times \C^4$.
Note that with respect to the product structure in~\eqref{eq:product-str-mfld}, we have
\begin{align}\label{eq:splitting-spinor-bundles}
    \pi_1^*\Sigma\T^2\otimes\pi_2^*\Sigma\sph^1\cong\Sigma\T^3=\T^3\times\C^2,
\end{align}
where
\begin{itemize}
    \item $\pi_1\colon\T^3\to\T^2$ and~$\pi_2\colon\T^3\to\sph^1$ are the projections,
    \item $\Sigma\T^2=\T^2\times \C^2$ is the trivial spinor bundle associated to the trivial spin structure of~$\T^2$,
    \item and~$\Sigma\sph^1=\sph^1\times\C$ is also the trivial spinor bundle over~$\sph^1$.
\end{itemize}
We refer to~\cite{Ginoux2009Dirac} and~\cite{Lawson1989Spin} for more information on spin geometry. 


\begin{rmk}
    There are two spin structures on the circle.
    Take the circle~$\sph^1(\frac{\ell}{2\pi})$ of perimeter~$\ell$ for example.
    The two spin structures are:
    \begin{itemize}
        \item the trivial spin structure~$\xi_0$, with spinors represented by~$\ell$-periodic complex-valued functions;
        \item the twisted spin structure~$\xi_1$, with spinors represented by~$\ell$-anti-periodic complex-valued functions. 
    \end{itemize}
    Let~$\theta\in [0,\ell]$ denote the angular coordinate for~$\sph^1(\frac{\ell}{2\pi})$, then the Dirac operator, in both cases, takes the form 
    \begin{align}
        \pd_{\sph^1}=-i\frac{\dd}{\dd\theta}.
    \end{align}
    However, the spectra for the Dirac operators with respect to different spin structures are different. 
    In the trivial case,
    \begin{align}
    \Spect(-i\frac{\dd}{\dd \theta};\xi_0)= \bZ \frac{2\pi}{\ell}, 
    \end{align}
    and for each~$\zeta\in\mathbb{Z}\frac{2\pi}{\ell}$, the eigenspace is generated by~$e^{i\zeta\theta}$. 
    Meanwhile for the twisted case, 
    \begin{align}
    \Spect(-i\frac{\dd}{\dd\theta};\xi_1)=(\frac{1}{2}+\mathbb{Z})\frac{2\pi}{\ell}
    \end{align}
    and for each~$\zeta\in(\frac{1}{2}+\mathbb{Z})\frac{2\pi}{\ell}$, the eigenspace for~$\zeta$ is again generated by~$e^{i\zeta\theta}$. 
    We refer to the note~\cite{Varilly2006Dirac} for further information. 
    Note that in either case, for any~$\zeta\in\Spect(-i\frac{\dd}{\dd \theta})$, the corresponding eigenspace is generated by~$e^{i\zeta\theta}$, and the eigenspinors have \emph{constant length}. 
\end{rmk}


We now clarify the spinors in this setting. 
The spinors on~$\R^{1+3}$ are~$\C^4$-valued functions. 
By the stationary ansatz, we are essentially dealing with spinors on~$\R^3$, via the relation 
\begin{align}
    \Sigma\R^{1+3}|_{\braces{t}\times \R^3}\cong \Sigma\R^3 \oplus \Sigma \R^3. 
\end{align}
Accordingly, we consider spinors in the bundle~$\Sigma\T^3\oplus\Sigma\T^3$, which corresponds to the up-down formulation of spinors as in~\cite{WZ2026Stationary}. 
Indeed, by writing
\begin{align}\label{eq:up-down}
    \psi= \begin{pmatrix}\psi^1\\ \psi^2 \\ \psi^3 \\ \psi^4\end{pmatrix}
    \equiv \begin{pmatrix} \psi_{\up} \\ \psi_{\down} \end{pmatrix}, 
\end{align}
the components~$\psi_{\up},\psi_{\down}$ are spinors on~$\Sigma \T^3$.
By~\eqref{eq:splitting-spinor-bundles}, we can further regard~$\psi_{\up},\psi_{\down}$ as a section of~$\pi_1^*\Sigma\T^2\otimes\pi_2^*\Sigma\sph^1$.
This motivates the following consideration: take~$\psi$ of the form
\begin{align}\label{eq:spinor-splitting}
    \psi= \begin{pmatrix}
        \psi_{\up} \\ \psi_{\down}
    \end{pmatrix}
    = 
    \begin{pmatrix}
        \varphi_{\up}\otimes \phi_{\up} \\ 
        \varphi_{\down}\otimes \phi_{\down}
    \end{pmatrix}
\end{align}
where~$\varphi_{\up},\varphi_{\down}$ are sections of~$\Sigma\T^2$ and~$\phi_{\up},\phi_{\down}$ are sections of~$\Sigma\sph^1$.
We hope to find nice candidates of~$\phi_{\up},\phi_{\down}$ such that spinors in the form of~\eqref{eq:spinor-splitting} can solve~\eqref{eq:NDE-spatial}.

\ 

As in~\cite{WZ2026Stationary} we write the massless and massive Dirac operators as 
\begin{align}
    \D=\sum_{k=1}^3 -i\gamma^0\gamma^k \p_k, & & \mbox{ respectively } & & 
    \pD=\sum_{k=1}^3 -i\gamma^0\gamma^k\p_k+m\gamma^0 =\D+m\gamma^0, 
\end{align}
and the intrinsic Dirac operator on~$\mathbb{T}^3$ equipped with the trivial spin structure is denoted as 
\begin{align}
    \pd_{\T^3}=\sum_{k=1}^3 -i\sigma^k \p_k = -i\sigma\cdot \nabla. 
\end{align}
Then the equation~\eqref{eq:NDE-spatial} acting on~$\psi$ of the form~\eqref{eq:spinor-splitting} is equivalent to 
\begin{align}
    \pd_{\T^3}\psi_{\down} & + m\psi_{\up} - a \psi_{\up} - (\p F)_{\up}=0,  \\
    \pd_{\T^3}\psi_{\up} & -m\psi_{\down} - a \psi_{\down} - (\p F)_{\down }=0. 
\end{align}
The intrinsic Dirac operator acts on tensor product spinors in the following manner: 
\begin{align}
    \pd_{\T^3} \psi_{\up} 
    =& \sum_{k=1}^3  -i\sigma^k\p_k (\varphi_{\up}\otimes \phi_{\up}) \\
    =& \sum_{k=1}^2 -i\sigma^k(\p_k\varphi_{\up})\otimes \phi_{\up}  + (-i\sigma^3)\varphi_{\up}\otimes \p_3\phi_{\up} \\
    =& \parenthesis{\sum_{k=1}^2 -i\sigma^k\p_k\varphi_{\up} }\otimes \phi_{\up}
       + \sigma^3\varphi_{\up}\otimes \parenthesis{-i\frac{\dd}{\dd x^3}\phi_{\up}} \\
    =& (\pd_{\T^2}\varphi_{\up})\otimes\phi_{\up}
        +\sigma^3\varphi_{\up}\otimes (\pd_{\sph^1}\phi_{\up})
\end{align}
and similarly 
\begin{align}
    \pd_{\T^3}\psi_{\down}
    = \pd_{\T^3} (\varphi_{\down}\otimes \phi_{\down}) 
    = (\pd_{\T^2} \varphi_{\down})\otimes \phi_{\down}
     + \sigma^3\varphi_{\down}\otimes (\pd_{\sph^1}\phi_{\down}). 
\end{align}
\begin{rmk}
    In geometric literature one finds the following convention:
    \begin{align}
        \pd_{\T^3} \psi_{\up} 
    = \pd_{\T^3} (\varphi_{\up}\otimes \phi_{\up}) 
    = \pd_{\T^2} \varphi_{\up}\otimes \phi_{\up}
     + \omega_{\T^2}^{\C}\cdot\varphi_{\up}\otimes \pd_{\sph^1}\phi_{\up},
    \end{align}
    where~$\omega_{\T^2}^{\C}$ stands for the complex volume element of~$\T^2$:
    \begin{align}
    \omega_{\T^2}^{\C}=i(-i\sigma^1)(-i\sigma^2).
    \end{align}
    By calculating explicitly using our choice of the Pauli matrices, we find that 
    \begin{align}
        \omega_{\T^2}^{\C}= \begin{pmatrix} 1 & 0 \\ 0 & -1 \end{pmatrix} = \sigma^3.
    \end{align}
    Thus the two formulas are consistent. 
\end{rmk}
Consequently, we can reformulate the equation~\eqref{eq:NDE-spatial} as a system for~$\varphi_{\up},\varphi_{\down}$ provided the other factors~$\phi_{\up},\phi_{\down}$ can be chosen appropriately. 
On the other hand, for the Lorentz scalar~$\bar{\psi}\psi$, which now takes the form 
\begin{align}
    \bar{\psi}\psi=\Abracket{\gamma^0\psi,\psi}
    =|\psi_{\up}|^2 - |\psi_{\down}|^2
    =|\varphi_{\up}|^2 |\phi_{\up}|^2
     - |\varphi_{\down}|^2 |\phi_{\down}|^2.
\end{align}
This suggests choosing both circle factors to have constant lengths. 

Here is the observation of splitting method: if we take~$\phi_{\up}=\phi_{\down}=e^{i \zeta x^3}$ with~$\zeta\in\mathbb{Z}\frac{2\pi}{\ell_3}$ which satisfies 
\begin{align}
    \pd_{\sph^1} e^{i\zeta x^3}= \zeta e^{i\zeta x^3},
\end{align}
then the spinor~$\psi$ of the form~\eqref{eq:spinor-splitting} satisfies 
\begin{align}
    \bar{\psi}\psi=\bar{\varphi}\varphi, 
    \qquad \mbox{ with } 
    \qquad 
    \varphi 
    = \begin{pmatrix}
    \varphi_{\up} \\ \varphi_{\down}
    \end{pmatrix}, 
    \quad \mbox{ and } \quad 
    \psi=\varphi\otimes e^{i\zeta x^3
    }.
\end{align}
If the nonlinearity~$F$ is also independent of the~$x^3$ variable, then 
\begin{align}
    \p F(\psi) 
    =& \p_\psi F((x^1,x^2), \psi(x)) 
    = \p_\psi F((x^1,x^2),\varphi(x^1,x^2) e^{i\zeta x^3} )\\
    =& e^{i\zeta x^3}\p_\psi F(x^1,x^2,\varphi(x^1,x^2))
    = e^{i\zeta x^3}\p F(\varphi)
\end{align}
since we have assumed that~$F$ is invariant under phase rotation from the very beginning. 
In this case the new spinor~$\varphi\in \Gamma(\Sigma \T^2\oplus \Sigma \T^2)$ satisfies the equation 
\begin{align}
    \begin{pmatrix}
        0 & \pd_{\T^2} + \zeta \omega_{\T^2}^{\C} \\ \pd_{\T^2}+ \zeta \omega_{\T^2}^{\C} & 0
    \end{pmatrix}\varphi 
    + m\gamma^0\varphi - a\varphi -\p F(\varphi)=0, 
\end{align}



Thus we can reduce~\eqref{eq:NDE-spatial} to a system in~$\varphi\in\Gamma(\Sigma\T^2\oplus \Sigma\T^2)$. 
We now rewrite this system in a convenient form. 
Note that 
\begin{align}
    \begin{pmatrix} 0 & \omega_{\T^2}^{\C} \\ \omega_{\T^2}^{\C} & 0 \end{pmatrix} 
    =
    \begin{pmatrix} 0 & \sigma^3 \\ \sigma^3 & 0 \end{pmatrix}
    = \gamma^0 \gamma^3  = -\gamma^3\gamma^0.
\end{align}
Thus if we fix a~$\zeta$ and denote
\begin{align}\label{eq:Dirac-operators}
    \D_{\T^2}
    \coloneqq \begin{pmatrix}  0 & \pd_{\T^2} \\ \pd_{\T^2} & 0 \end{pmatrix}= -i\sum_{k=1}^2 \gamma^0\gamma^k\p_k, 
    & & \mbox{ and } & & 
    \pD_{\T^2} \coloneqq \D_{\T^2}+\zeta\gamma^0\gamma^3+m\gamma^0, 
\end{align}
then the equation in~$\varphi$ takes the form
\begin{align}\label{eq:NDE-T2}
    \pD_{\T^2}\varphi- a\varphi -\p F(\varphi)=0. 
\end{align}
This is an elliptic system on the two-dimensional torus~$\T^2= \faktor{\R^2}{\Gamma(\ell_1,\ell_2)}$. 


\section{Spectral properties of the Dirac operators}\label{sect:spectral}

On a general closed spin manifold~$M$, the qualitative spectral properties of the operators 
\begin{align}
    \pd_M + f(x)\omega^{\C}_M, \qquad \mbox{ with } f\in C^\infty(M, \R_+)
\end{align}
are investigated in~\cite{SireXu2023variational}. 
Here we need more quantitative spectral properties. 
In our setting the coefficients are constant, and the torus Dirac spectra can be computed explicitly. 
Observe from the Clifford relations that
    \begin{align}
        \D_{\T^2}\parenthesis{\zeta\gamma^0\gamma^3+m\gamma^0}
        = -\parenthesis{\zeta\gamma^0\gamma^3+m\gamma^0} \D_{\T^2}, 
        & & 
        \mbox{ and } 
        & & 
        \parenthesis{\zeta\gamma^0\gamma^3+m\gamma^0}^2
        =\parenthesis{\zeta^2+m^2}I_4.
    \end{align}
Thus we have
\begin{align}
    \pD_{\T^2}^2
    =\parenthesis{ \D_{\T^2} + \zeta \gamma^0\gamma^3 + m\gamma^0}^2 
    =\D_{\T^2}^2 + \parenthesis{\zeta^2+m^2}I_4.
\end{align}
where~$I_4$ denotes the~$4\times 4$ identity matrix.
We first determine the spectrum of massless Dirac operator~$\D_{\T^2}$. 

Recall that the eigenvalues of~$\pd_{\T^2}$ associated to the trivial spin structure on the torus~$\T^2=\faktor{\R^2}{\Gamma(\ell_1,\ell_2)}$ are computed in~\cite{Friedrich1984zur} and~\cite{Ginoux2009Dirac}:  
\begin{align}
    \Spect(\pd_{\T^2}) =
    \braces{ \pm 2\pi |\zeta^*| \mid \zeta^*\in\Gamma(\ell_1,\ell_2)^*},
\end{align}
and the complex multiplicity of~$0$ is~$2$ (with eigenspinors given by the two linearly independent constant spinors, say~$(1,0)^T$ and~$(0,1)^T$), while the complex multiplicity of the nonzero eigenvalue~$2\pi|\zeta^*|$ equals the number of elements in~$\Gamma(\ell_1,\ell_2)^*$ whose length equals~$|\zeta^*|$.
This is a symmetric set with respect to the origin in the real line, whether or not multiplicities are counted. 
We enumerate the eigenvalues of~$\pd_{\T^2}$ in a non-decreasing order with multiplicities counted: 
\begin{align}\label{eq:spectrum-pd}
        \cdots \leq \lambda_{-k}\leq  \cdots \leq \lambda_{-1} <0=\lambda_{0,1}=\lambda_{0,2}<\lambda_1\leq \cdots \leq \lambda_k \cdots  
    \end{align}
with~$\lambda_{-k}=-\lambda_k$. 
Let~$\braces{\eta_k\mid k\in\mathbb{Z},k\neq 0} \cup \braces{\eta_{0,1},\eta_{0,2}}$ be a complete orthonormal basis of~$L^2(\T^2,\C^2)$ consisting of eigenspinors:
\begin{align}
    \pd_{\T^2}\eta_k = \lambda_k \eta_k, \qquad \forall k\neq 0
\end{align}
\begin{align}
    \pd_{\T^2}\eta_{0,j}=\lambda_{0,j}\eta_{0,j}=0, \quad j=1,2. 
\end{align}
Then, similar to the discussion in~\cite{WZ2026Stationary}, we have 
\begin{itemize}
    \item[(1)] $\Spect(\pd_{\T^2})\subset\Spect(\D_{\T^2})$, and for each~$\lambda_k\in\Spect(\pd_{\T^2})\setminus \braces{0}$,
    \begin{align}
        \begin{pmatrix} \eta_k \\ \eta_k \end{pmatrix}, 
        \begin{pmatrix} \eta_{-k} \\ -\eta_{-k} \end{pmatrix}
        \in \Eigen(\D_{\T^2};\lambda_k)
    \end{align}
    and the above two spinors are linearly independent (but not normalized).
    Meanwhile the~$0$-eigenspinors of~$\D_{\T^2}$ are given by the constant sections 
    \begin{align}
        \psi_{0,1}=\begin{pmatrix} 1 \\ 0 \\ 0 \\ 0 \end{pmatrix}, \quad 
        \psi_{0,2}=\begin{pmatrix} 0 \\ 1 \\ 0 \\ 0 \end{pmatrix}, \quad 
        \psi_{0,3}=\begin{pmatrix} 0 \\ 0 \\ 1 \\ 0 \end{pmatrix}, \quad 
        \psi_{0,4}=\begin{pmatrix} 0 \\ 0 \\ 0 \\ 1 \end{pmatrix}. 
    \end{align}

    \item[(2)] $\Spect(\D_{\T^2})\subset\Spect(\pd_{\T^2})$, and if~$\D_{\T^2}\psi=\lambda\psi$, then
    \begin{align}
        \pd_{\T^2}(\psi_{\up}+\psi_{\down}) = \lambda (\psi_{\up}+\psi_{\down}), & & 
        \pd_{\T^2}(\psi_{\up}-\psi_{\down}) =-\lambda (\psi_{\up}-\psi_{\down}).
    \end{align}
    Together with the construction in (1) we see that both~$\lambda$ and~$-\lambda$ are in~$\Spect(\pd_{\T^2})$. 
\end{itemize}
Moreover, the operator~$\D_{\T^2}$ is unitarily equivalent to~$\pd_{\T^2}\oplus (-\pd_{\T^2})$ under the transform 
\begin{align}
    \begin{pmatrix}\psi_{\up} \\ \psi_{\down} \end{pmatrix}
    \mapsto 
    \frac{1}{\sqrt{2}}
    \begin{pmatrix} \psi_{\up}+\psi_{\down} \\ \psi_{\up}-\psi_{\down} \end{pmatrix}. 
\end{align}
Recalling~$\Spect(\pd_{T^2})$ is symmetric in~$\R$, we thus have 
\begin{align}
    \Eigen(\D_{\T^2};\lambda) \cong \Eigen(\pd_{\T^2};\lambda)\oplus\Eigen(\pd_{\T^2};-\lambda), & & 
    \Mtp(\D_{\T^2};\lambda)= 2\Mtp(\pd_{\T^2};\lambda).
\end{align}
This proves the first half of the following lemma. 

\begin{lemma}\label{lemma:spectrum}
    Let~$\mathbb{T}^2=\faktor{\R^2}{\Gamma(\ell_1,\ell_2)}$ be equipped with the trivial spin structure, and let~$\D_{\T^2}, \pD_{\T^2}$ be defined as in~\eqref{eq:Dirac-operators}. 
    \begin{itemize}
        \item[(i)] The spectrum of~$\D_{\T^2}$ is 
        \begin{align}
            \Spect(\D_{\T^2})= \Spect(\pd_{\T^2})
            =\braces{\pm 2\pi |\zeta^*| \mid \zeta^*\in \Gamma(\ell_1,\ell_2)^*\subset\R^2}
        \end{align}
        and for each~$\lambda\in \Spect(\pd_{\T^2})$,~$\Mtp(\D_{\T^2};\lambda)= 2\Mtp(\pd_{\T^2};\lambda)$.

        \item[(ii)] The spectrum of~$\pD_{\T^2}$ is 
            \begin{align}
                \Spect(\pD_{\T^2})=\braces{ \pm \sqrt{\zeta^2+\lambda^2+m^2} \mid \lambda \in \Spect(\D_{\T^2})} 
            \end{align}
            with multiplicities given by
            \begin{align}
                \Mtp(\pD_{\T^2};\sqrt{\zeta^2+\lambda^2+m^2})
                =\Mtp(\pD_{\T^2};-\sqrt{\zeta^2+\lambda^2+m^2})
                =\Mtp(\D_{\T^2}; \lambda),
            \end{align}
            for any~$\lambda\in \Spect(\D_{\T^2})\setminus \braces{0}$, while 
            \begin{align}
                \Mtp(\pD_{\T^2};\sqrt{\zeta^2+m^2})
                =\Mtp(\pD_{\T^2};-\sqrt{\zeta^2+m^2})=2.
            \end{align}
    \end{itemize}
    
\end{lemma}

\begin{proof}
    It remains to determine the spectrum of~$\pD_{\T^2}$.

    First note that the operator~$\D_{\T^2}^2$ is an elliptic operator of Laplacian type. 
    It is also self-adjoint on a closed manifold. 
    Every eigenspinor of~$\D_{\T^2}$ is automatically an eigenspinor of~$\D_{\T^2}^2$.
    Since such eigenspinors form a complete basis for the~$L^2$ space of sections of the bundle~$\Sigma\T^2\oplus \Sigma\T^2$, we see that
    \begin{align}
        \Spect(\D_{\T^2}^2)=\braces{\lambda^2 \mid \lambda\in\Spect(\D_{\T^2})}
    \end{align}
    and for~$\lambda\neq 0$, 
    \begin{align}\label{eq:Eigen-lambda2}
        \Eigen(\D_{\T^2}^2;\lambda^2)
        =\Eigen(\D_{\T^2};\lambda)
        \oplus
        \Eigen(\D_{\T^2};-\lambda),
    \end{align}
    \begin{align}
        \Mtp(\D_{\T^2}^2;\lambda^2)
        = \Mtp(\D_{\T^2};\lambda) 
        + \Mtp(\D_{\T^2};-\lambda)
        =2\Mtp(\D_{\T^2};\lambda), 
    \end{align}
    while~$\Mtp(\D_{\T^2}^2;0)=4=\Mtp(\D_{\T^2};0)$.

    As a self-adjoint elliptic operator on a closed manifold, the operator~$\pD_{\T^2}$ also has only eigenvalues of finite multiplicities. 
    Suppose~$\rho\in\Spect(\pD_{\T^2})$ with~$\pD_{\T^2}\psi=\rho\psi$.
    Then 
    \begin{align}
        \rho^2\psi=\pD_{\T^2}^2\psi 
        = \D_{\T^2}^2\psi + (\zeta^2+m^2)\psi. 
    \end{align}
    Thus~$\D_{\T^2}^2\psi= (\rho^2-\zeta^2-m^2)\psi=\lambda^2\psi$ for some~$\lambda\in\Spect(\D_{\T^2})$. 
    It follows that
    \begin{align}
        \rho\in \braces{\pm \sqrt{\zeta^2+\lambda^2+m^2}}. 
    \end{align}
    Next we show that both of~$\pm \sqrt{\zeta^2+\lambda^2+m^2} $ appear in the spectrum of~$\pD_{\T^2}$, with the same multiplicity as~$\lambda\in\Spect(\D_{\T^2})$.

    
    For each~$\lambda>0$, consider the linear map
    \begin{align}
        \frac{\zeta\gamma^0\gamma^3+m\gamma^0}{\sqrt{\zeta^2+m^2}}\colon \Eigen(\D_{\T^2};\lambda)\to \Eigen(\D_{\T^2};-\lambda).
    \end{align}
    It is well-defined: if~$\D_{\T^2}\psi=\lambda\psi$, then 
    \begin{align}
        \D_{\T^2}\parenthesis{\frac{1}{\sqrt{\zeta^2+m^2}}(\zeta\gamma^0\gamma^3+m\gamma^0)\psi}
        =&-\frac{1}{\sqrt{\zeta^2+m^2}}\parenthesis{\zeta\gamma^0\gamma^3+m\gamma^0}
        \D_{\T^2}\psi \\
        =&-\lambda \parenthesis{\frac{1}{\sqrt{\zeta^2+m^2}}(\zeta\gamma^0\gamma^3+m\gamma^0)\psi}. 
    \end{align}
    It is also invertible since 
    \begin{align}
        \parenthesis{ \frac{\zeta\gamma^0\gamma^3+m\gamma^0}{\sqrt{\zeta^2+m^2}} }^2=I_4.
    \end{align}
    Let~$d\equiv \Mtp(\D_{\T^2};\lambda)$ and choose an orthonormal basis~$\psi_1,\dots, \psi_d$ of~$\Eigen(\D_{\T^2};\lambda)$. 
    Set
    \begin{align}
        \phi_j \coloneqq \frac{1}{\sqrt{\zeta^2+m^2}} \parenthesis{\zeta\gamma^0\gamma^3+m\gamma^0}\psi_j,\quad 1\leq j\leq d. 
    \end{align}
    Then~$\phi_1,\dots,\phi_d$ form a basis of~$\Eigen(\D_{\T^2};-\lambda)$. 
    Moreover,
    \begin{align}
        \parenthesis{\zeta\gamma^0\gamma^3+m\gamma^0}\psi_j= \sqrt{\zeta^2+m^2}\phi_j, 
        & & 
        \parenthesis{\zeta\gamma^0\gamma^3+m\gamma^0}\phi_j = \sqrt{\zeta^2+m^2}\psi_j.
    \end{align}
    Consequently,
    \begin{align}
        \pD_{\T^2}\psi_j = \lambda \psi_j+\sqrt{\zeta^2+m^2}\phi_j,
        & & 
        \pD_{\T^2}\phi_j = \sqrt{\zeta^2+m^2}\psi_j-\lambda \phi_j.
    \end{align}
    Hence the two-dimensional space~$\Span_{\C}\{\psi_j,\phi_j\}$ is invariant under~$\pD_{\T^2}$,
    on which~$\pD_{\T^2}$ acts via a matrix
    \begin{align}
        \begin{pmatrix}
            \lambda & \sqrt{\zeta^2+m^2}\\
            \sqrt{\zeta^2+m^2} & -\lambda
        \end{pmatrix}.
    \end{align}
    Its characteristic polynomial is~$p(\mu)=\mu^2-\lambda^2-\zeta^2-m^2$, and thus its eigenvalues are
    \begin{align}
        \mu=\pm\sqrt{\lambda^2+\zeta^2+m^2}.
    \end{align}

    Recall~$\Eigen(\D_{\T^2}^2;\lambda^2)$ is decomposed as in~\eqref{eq:Eigen-lambda2}.
    The above construction decomposes this space into~$d$ two-dimensional invariant subspaces. Hence each of the eigenvalues~$\pm\sqrt{\lambda^2+\zeta^2+m^2}$ has multiplicity~$d$.  
    Hence
    \begin{align}
        \Mtp\parenthesis{
            \pD_{\T^2};
            \sqrt{\lambda^2+\zeta^2+m^2}}
        =\Mtp\parenthesis{
            \pD_{\T^2};
            -\sqrt{\lambda^2+\zeta^2+m^2}}
        =\Mtp(\D_{\T^2};\lambda).
    \end{align}

    As for the eigenspaces for~$\pm\sqrt{\zeta^2+m^2}$, we simply note that they must arise from~$\Eigen(\D_{\T^2};0)$ and~$\pD_{\T^2}$ acting on constant spinors reduces to the matrix
    \begin{align}
        \mathscr{D}\equiv
        \begin{pmatrix}
            m & 0 & \zeta & 0 \\
            0 & m & 0& -\zeta \\
            \zeta & 0 & -m & 0 \\
            0 & -\zeta & 0 & -m
        \end{pmatrix}. 
    \end{align}
    This matrix has eigenvalues~$\pm\sqrt{\zeta^2+m^2}$ with eigenspaces 
    \begin{align}
        \Eigen(\mathscr{D};\sqrt{\zeta^2+m^2})
        =\Span\braces{ \begin{pmatrix} 1 \\ 0 \\ \frac{\zeta}{ m+\sqrt{\zeta^2+m^2} } \\ 0  \end{pmatrix}, 
        \begin{pmatrix} 0 \\ 1 \\ 0 \\ \frac{-\zeta}{m+\sqrt{\zeta^2+m^2}} \end{pmatrix} }, 
    \end{align}
    \begin{align}
        \Eigen(\mathscr{D};-\sqrt{\zeta^2+m^2})
        =\Span\braces{ \begin{pmatrix} \frac{-\zeta}{m+\sqrt{\zeta^2+m^2}} \\ 0 \\ 1 \\ 0  \end{pmatrix}, 
        \begin{pmatrix} 0 \\ \frac{\zeta}{m+\sqrt{\zeta^2+m^2}} \\ 0 \\ 1 \end{pmatrix}}. 
    \end{align}
    Note that these constant eigenspinors reduce to the~$\psi_{0,j}$'s when~$\zeta=0$. 
    This proves the second part of the lemma. 
\end{proof}

\section{Null-cone separation and a resolvent estimate}\label{sect:key lemma}

The space~$\Eigen(\pD_{\T^2}, \sqrt{\zeta^2+m^2})$ consists of constant spinors on $\mathbb{T}^2$, but such spinors no longer lie in the $(+1)$-eigenspace of $\gamma^0$, a property used in the contradiction argument in~\cite{WZ2026Stationary}.
Nevertheless, they cannot lie in the Lorentz null cone.  
Indeed, any~$e\in\Eigen(\pD_{\T^2};\sqrt{\zeta^2+m^2})$ satisfies
\begin{align}
    \bar{e}e=\frac{m}{\sqrt{m^2+\zeta^2}}|e|^2.
\end{align}
which is nonzero unless~$e=0$. 
In particular,~$\mathcal{N}\cap \Eigen(\pD_{\T^2};\sqrt{\zeta^2+m^2})=\braces{0}$.  

Let~$\proj_{m,\zeta}$ be the~$L^2$-orthogonal projection to~$\Eigen(\pD_{\T^2};\sqrt{m^2+\zeta^2})$ and~$\proj_{m,\zeta}^\perp = \id -\proj_{m,\zeta}$ the complementary projection.
The following lemma gives a quantitative separation of the null cone from the lowest positive eigenspace~$\Eigen(\pD_{\T^2};\sqrt{\zeta^2+m^2})$. 
\begin{lemma}\label{lemma:null-cone-separation}
    For any~$\nu\in (1,2)$, there exist~$\kappa_1>0$ and~$\kappa_2>0$ such that
    \begin{align}\label{eq:kappa}
        \|\varphi\|_{L^4}=1,\quad \|\bar{\varphi} \varphi\|_{L^\nu}\leq\kappa_1 \quad\Longrightarrow\quad \|\proj_{m,\zeta}^\perp \varphi\|_{L^4}\geq\kappa_2.
    \end{align}
\end{lemma}
\begin{proof}
    The corresponding statement for~$\nu\in(1,\frac{3}{2})$ was proved in~\cite{WZ2026Stationary}; but the argument works also for~$\nu\in (1,2)$.  
    Since these constants are crucial for us, we give the detailed proof here. 
    
    To see~\eqref{eq:kappa}, we argue by contradiction. 
    Suppose  there exists a sequence~$(\varphi_k)$ such that
    \begin{align}
        \|\varphi_k\|_{L^4}=1, \quad \|\bar{\varphi}_k\varphi_k\|_{L^\nu} \to 0,
        \quad \|\proj_{m,\zeta}^\perp \varphi_k\|_{L^4}\to 0, \qquad \mbox{ as } \; k\to+\infty.
    \end{align}
    We observe that ~$\|\proj_{m,\zeta}\varphi_k\|_{L^4}\leq C\|\varphi_k\|_{L^4}=C$, with~$C$ depending on~$\T^2$ and the finite-dimensional space~$\Eigen(\pD_{\T^2};\sqrt{\zeta^2+m^2})$. 
    Since any two norms on a finite-dimensional space are equivalent, the sequence~$(\proj_{m,\zeta}\varphi_k)$ is bounded in~$\Eigen(\pD_{\T^2};\sqrt{m^2+\zeta^2})\cong\C^2$. 
    Therefore, upon passing to a subsequence if necessary, the sequence~$\proj_{m,\zeta}\varphi_k$ converges to some~$\varphi_{*}$.
    Together with the convergence~$\|\proj_{m,\zeta}^\perp \varphi_k\|_{L^4}\to0$, we see that~$\varphi_k$ converges in~$L^4$ to~$\varphi_*$ with~$\|\varphi_{*}\|_{L^4}=1$, and~$\bar{\varphi}_k\varphi_k$ converges in~$L^2$ to~$\bar{\varphi}_*\varphi_*$. 
    Since~$\nu<2$, we see that~$\bar{\varphi}_k\varphi_k$ also converges in~$L^\nu$ to~$\bar{\varphi}_*\varphi_*$, hence~$\bar{\varphi}_*\varphi_*=0$ a.e. 
    This is impossible, since~$\varphi_*\in\Eigen(\pD_{\T^2};\sqrt{m^2+\zeta^2})$ and~$\bar{\varphi}_*\varphi_*= \frac{m}{\sqrt{\zeta^2 + m^2}}|\varphi_*|^2$.
    
\end{proof}
      We remark that the above statement actually holds for any~$\nu\geq 1$ and there is no need for the requirement~$\nu<2$.
      This follows by a H\"older inequality argument. 
      Indeed, there holds an even more general principle. 
      We state it here for future reference. 

\begin{lemma}[An abstract separation lemma]\label{lemma:abstract-separation}
    Let~$X$ be a normed linear space and let~$E\subset X$ be a finite-dimensional subspace.
    Let~$(Z,d_Z)$ be a metric space with a distinguished point~$0_Z\in Z$, and let
    \begin{align}
        \Phi\colon X\to Z
    \end{align}
    be continuous.
    Assume that
    \begin{align}\label{eq:trivial-intersection}
        E\cap \Phi^{-1}(0_Z)=\braces{0}.
    \end{align}
    Suppose moreover that~$\rho\colon X\to[0,+\infty)$ has the property that, for every sequence~$(u_n)\subset X$,
    \begin{align}\label{eq:rho-Phi}
        \rho(u_n)\to0
        \qquad\Longrightarrow\qquad
        d_Z(\Phi(u_n),0_Z)\to0.
    \end{align}
    Then there exist constants~$\kappa_1,\kappa_2>0$ such that
    \begin{align}\label{eq:abstract-separation}
        \|u\|_X=1,
        \qquad
        \rho(u)\leq\kappa_1
        \qquad\Longrightarrow\qquad
        \dist_X(u,E)\geq\kappa_2.
    \end{align}
\end{lemma}

\begin{proof}
    Suppose not, then there would exist a sequence~$(u_n)\subset X$ such that
    \begin{align}
        \|u_n\|_X=1,
        \qquad
        \rho(u_n)\leq\frac{1}{n},
        \qquad
        \dist_X(u_n,E)\leq\frac{1}{n}.
    \end{align}
    Choose~$e_n\in E$ such that~$\|u_n-e_n\|_X\leq \frac{2}{n}$. 
    In particular,~$\|e_n\|_X\to1$. 
    Since~$E$ is finite dimensional, after passing to a subsequence if necessary, there exists~$e_*\in E$ such that~$e_n\to e_*$ in~$X$. 
    Consequently,
    \begin{align}
        u_n\to e_*
        \qquad\mbox{ in }X,
        \qquad
        \|e_*\|_X=1.
    \end{align}
    By the continuity of~$\Phi$,
    \begin{align}
        d_Z(\Phi(u_n),\Phi(e_*))\to0.
    \end{align}
    On the other hand,~$\rho(u_n)\to0$, and hence~\eqref{eq:rho-Phi} gives~$ d_Z(\Phi(u_n),0_Z)\to0$. 
    By uniqueness of limits in the metric space~$Z$, we must have~$\Phi(e_*)=0_Z$. 
    Thus~$e_*\in E\cap\Phi^{-1}(0_Z)=\braces{0}$, contradicting~$\|e_*\|_X=1$.
\end{proof}

\begin{rmk}
    To apply the above general principle to our case, we let
    \begin{align}
        X=L^p(\T^2,\C^4),
        \qquad
        \Phi(\varphi)=\bar{\varphi}\varphi,
        \qquad
        \rho(\varphi)=\|\bar{\varphi}\varphi\|_{L^q}.
    \end{align}
    with~$1\leq p,q<+\infty$. 
    Set
    \begin{align}
        r\coloneqq \min\braces{q,\frac{p}{2}}.
    \end{align}
    If~$r\geq1$, we take~$Z=L^r(\T^2)$ with its usual norm metric, while if~$0<r<1$, we equip~$L^r(\T^2)$ with the metric
    \begin{align}
        d_r(f,g)\coloneqq \int_{\T^2}|f-g|^r\dv_{\T^2}.
    \end{align}
    Since~$r\leq q$ and~$\T^2$ has finite volume,
    \begin{align}
        \|\bar{\varphi}_n\varphi_n\|_{L^q}\to0
        \qquad\Longrightarrow\qquad
        \bar{\varphi}_n\varphi_n\to0
        \quad\mbox{ in }Z.
    \end{align}
    Moreover, since~$2r\leq p$, the map
    \begin{align}
        \Phi\colon L^p(\T^2,\C^4)\to Z,
        \qquad
        \varphi\mapsto\bar{\varphi}\varphi,
    \end{align}
    is continuous.
    Hence, for every finite-dimensional subspace~$E\subset L^p(\T^2,\C^4)$ satisfying
    \begin{align}
        E\cap\braces{\varphi\mid \bar{\varphi}\varphi=0\;\mbox{ a.e.}}=\braces{0},
    \end{align}
    there exist~$\kappa_1,\kappa_2>0$ such that
    \begin{align}
        \|\varphi\|_{L^p}=1,
        \qquad
        \|\bar{\varphi}\varphi\|_{L^q}\leq\kappa_1
        \qquad\Longrightarrow\qquad
        \dist_{L^p}(\varphi,E)\geq\kappa_2.
    \end{align}
    In our case,~$E=\Eigen(\pD_{\T^2};\sqrt{\zeta^2+m^2})$, and~$\dist_{L^p}(\varphi,E)\leq \|\proj_{m,\zeta}^\perp \varphi\|_{L^p}$, so the general lemma applies. 
\end{rmk}
Note that the above compactness argument is qualitative and does not provide explicit values of~$\kappa_1$ and~$\kappa_2$.


\ 

      We also need a resolvent estimate for~$\pD_{\T^2}-a $, for any~$a\in [0,\sqrt{\zeta^2+m^2}]$, on the orthogonal complement of the lowest positive eigenspace~$\Eigen(\pD_{\T^2};\sqrt{\zeta^2+m^2})$.

\begin{lemma}\label{lemma:invertibility-orthogonal}
    There exists a constant~$C=C(m,\zeta,\mathbb{T}^2)$ such that for any~$a\in [0,\sqrt{\zeta^2+m^2}]$, 
    \begin{align}
        \|\proj_{m,\zeta}^\perp \varphi\|_{L^4}\leq C\|\proj_{m,\zeta}^\perp(\pD_{\T^2}-a)\varphi\|_{L^{\frac{4}{3}}}.
    \end{align}
\end{lemma}

\begin{proof}
    Suppose not, then there exist sequences~$(a_k)\subset [0,\sqrt{\zeta^2+m^2}]$ and~$(\varphi_k)\subset W^{1,\frac{4}{3}}$ such that
        \begin{align}
            \|\proj_{m,\zeta}^\perp\varphi_k\|_{L^4}=1, & &
            \|\proj_{m,\zeta}^\perp(\pD_{\T^2}-a_k)\varphi_k\|_{L^{\frac{4}{3}}}\to 0.
        \end{align}
        We may assume~$a_k\to a_*\in [0,\sqrt{\zeta^2+m^2}]$ by Heine--Borel theorem. 
        Note that for any~$k$,
        \begin{align}
            \|\proj_{m,\zeta}^\perp \varphi_k \|_{W^{1,\frac{4}{3}}}
            \leq& C \parenthesis{ \| \proj_{m,\zeta}^\perp \varphi_k\|_{L^{\frac{4}{3}}} + \|\pD_{\T^2}\proj_{m,\zeta}^\perp \varphi_k\|_{L^{\frac{4}{3}}} } \\
            \leq& C \parenthesis{
            \| \proj_{m,\zeta}^\perp \varphi_k\|_{L^{\frac{4}{3}}}
            + \|(\pD_{\T^2}-a_k)\proj_{m,\zeta}^\perp \varphi_k\|_{L^{\frac{4}{3}}} 
            + a_k\|\proj_{m,\zeta}^\perp \varphi_k\|_{L^{\frac{4}{3}}} } \\
            \leq& C \parenthesis{ (\sqrt{m^2+\zeta^2}+1)\| \proj_{m,\zeta}^\perp \varphi_k\|_{L^{\frac{4}{3}}} + \|\proj_{m,\zeta}^\perp (\pD_{\T^2}-a_k)\varphi_k\|_{L^{\frac{4}{3}}} },
        \end{align}
        where~$C$ is independent of~$k$ and we have used the commutativity of~$\proj_{m,\zeta}^\perp$ with~$(\pD_{\T^2}-a_k)$ in the last inequality.
        Hence the sequence~$(\proj_{m,\zeta}^\perp \varphi_k)$ is also uniformly bounded in~$W^{1,\frac{4}{3}}$ and we may assume
        \begin{align}
            \proj_{m,\zeta}^\perp\varphi_k \to \varphi_*=\proj_{m,\zeta}^\perp \varphi_*\quad  \mbox{ in } \;  L^{\frac{4}{3}}.
        \end{align}
        Moreover, from
        \begin{align}
            \pD_{\T^2}(\proj_{m,\zeta}^\perp\varphi_k)= \proj_{m,\zeta}^\perp (\pD_{\T^2}-a_k)\varphi_k+ a_k \proj_{m,\zeta}^\perp\varphi_k \xrightarrow{L^{\frac{4}{3}}} 0+ a_*\varphi_*
        \end{align}
        it follows that
        \begin{align}
            \pD_{\T^2}\varphi_*=a_*\varphi_*.
        \end{align}
        Since~$\varphi_* =\proj_{m,\zeta}^\perp\varphi_*\perp \Eigen(\pD_{\T^2};\sqrt{m^2+\zeta^2})$ while~$\dist(a_*, \Spect(\pD_{\T^2})\setminus\braces{\sqrt{m^2+\zeta^2}})>0$, we must have~$\varphi_*=0$.
        This in turn implies that
        \begin{align}
            \|\proj_{m,\zeta}^\perp \varphi_k\|_{W^{1,\frac{4}{3}}} \to 0
        \end{align}
        as~$k\to+\infty$, contradicting~$\|\proj_{m,\zeta}^\perp\varphi_k\|_{L^4}=1$ for any~$k$.
\end{proof}

Another explanation of this resolvent estimate will be given in Remark~\ref{rmk:resolvent-estimate}, using Fourier modes, and we can make the constant~$C$ more explicit. 



\section{The variational framework}\label{sect:variational framework}

In this section we build up the working space~$H^{\frac{1}{2}}(\T^2;\C^4)$ and realize the equation~\eqref{eq:NDE-T2} as the Euler-Lagrange equation of a~$C^1$ functional, so that we can obtain nontrivial solutions via minimax method in the sequel. 

As explained above, we consider the bundle~$\Sigma\T^2\oplus\Sigma\T^2\to\T^2$, and work with sections of regularity~$H^{1/2}$.  
This can be defined using the operator~$\pD_{\T^2}$, the resulting spaces agree with the usual fractional Sobolev spaces, see e.g.~\cite{Ammann2003Habil, SireXu2023variational} for more details. 
The construction is analogous to that in~\cite{WZ2026Stationary}; we include the details needed below.  


To the self-adjoint elliptic operator~$\pD_{\T^2}$ we associate the spectral measure
\begin{align}
    \tau\colon \mathcal{B}(\R)\to \proj(L^2(\T^2;\C^4))
\end{align}
Then we have 
\begin{align}
    \pD_{\T^2} = \int_{\R} \lambda \dd\tau(\lambda)
\end{align}
and the absolute value operator of~$\pD_{\T^2}$ is given formally by 
\begin{align}
    |\pD_{\T^2}|=\int_\R |\lambda| \dd\tau(\lambda).
\end{align}
For any~$s\geq 0$, we can define the fractional powers~$|\pD_{\T^2}|^s$ by 
\begin{align}
    |\pD_{\T^2}|^s = \int_{\R} |\lambda|^s \dd\tau(\lambda)
\end{align}
whose domain is the space
\begin{align}
    \Dom(|\pD_{\T^2}|^s)
    =&\braces{ \varphi\in L^2(\T^2,\C^4) \mid \int_\R |\lambda|^{2s} \dd\Abracket{\tau(\lambda)\varphi,\varphi } <+\infty}  \\
    =&H^{s}(\mathbb{T}^2;\C^4) 
    =H^s(\T^2;\Sigma\T^2\oplus\Sigma\T^2).
\end{align}
The space~$H^{-s}(\T^2,\C^4)$ denotes the dual space of~$H^{s}(\T^2,\C^4)$ as usual. 

\ 

In our situation we need to consider~$s=\frac{1}{2}$, which is the most natural space to deal with Dirac equations. 
The working space is 
\begin{align}
    H^{\frac{1}{2}}(\T^2,\C^4)
    =\braces{\varphi\in L^2(\T^2,\C^4) \; \mid \; \|\varphi\|_{H^{\frac{1}{2}}} \equiv \parenthesis{ \|\varphi\|_{L^2}^2 + \| |\pD_{\T^2}|^{\frac{1}{2}} \varphi \|_{L^2}^2 }^{\frac{1}{2}} <+\infty  }. 
\end{align}
Since~$m>0$, the following norm 
\begin{align}
    \tnorm{\varphi}_{H^{\frac{1}{2}}(\T^2,\C^4)}\coloneqq \| |\pD_{\T^2}|^{\frac{1}{2}}\varphi\|_{L^2(\T^2,\C^4)}
\end{align}
is equivalent to the standard norm~$\|\cdot\|_{H^{\frac{1}{2}}}$. 

As usual we decompose the Hilbert space~$H^{\frac{1}{2}}$ as the direct sum of the subspaces spanned by eigenspinors of positive respectively negative eigenvalues
\begin{align}
    H^{\frac{1}{2}}(\mathbb{T}^2,\C^4)= H^{\frac{1}{2},+}\oplus H^{\frac{1}{2},-}, 
\end{align}
with 
\begin{align}
    H^{\frac{1}{2},+}\coloneqq \overline{ \bigoplus_{k>0} \Eigen(\pD_{\T^2};\mu_k) }^{H^{\frac{1}{2}}}, \quad 
    H^{\frac{1}{2},-}\coloneqq \overline{ \bigoplus_{k<0} \Eigen(\pD_{\T^2};\mu_k) }^{H^{\frac{1}{2}}} . 
\end{align}
The orthogonal projections onto these subspaces are denoted by~$P^{\pm}$. 
With these notations, we have 
\begin{align}
    \varphi=P^+\varphi + P^-\varphi \equiv \varphi^+ + \varphi^- \in H^{\frac{1}{2},+}\oplus H^{\frac{1}{2},-}=H^{\frac{1}{2}}, 
\end{align}
\begin{align}
    \pD_{\T^2}=&\pD_{\T^2}(P^+ + P^-) = \pD_{\T^2} P^+ + \pD_{\T^2} P^-,  \\
    |\pD_{\T^2}|=& \pD_{\T^2} P^+ - \pD_{\T^2} P^-. 
\end{align}
Note that~$\pD_{\T^2}$ is a first-order elliptic operator, and for a spinor~$\varphi\in H^{\frac{1}{2}}(\T^2,\C^4)$, 
\begin{align}
    \pD_{\T^2}\varphi = \pD_{\T^2}P^+\varphi + \pD_{\T^2}P^-\varphi 
    =|\pD_{\T^2}|P^+\varphi - |\pD_{\T^2}|P^-\varphi 
\end{align}
is only an element in~$H^{-\frac{1}{2}}$.
Thus~$\Abracket{\pD_{\T^2}\varphi,\varphi}$ cannot be understood as a~$L^2$ function. 
However, as usually done in analysis, we can make sense of it via the duality pairing~$\Abracket{\pD_{\T^2}\varphi,\varphi}_{H^{-1/2}\times H^{1/2}}\in \R$, and write 
\begin{align}
    \int_{\T^2}\Abracket{\pD_{\T^2}\varphi,\varphi}\dv_{\T^2}
    \coloneqq 
    \Abracket{\pD_{\T^2}\varphi,\varphi}_{H^{-1/2}\times H^{1/2}}\in \R. 
\end{align}
In terms of the above decompositions, we have 
\begin{align}
    \int_{\T^2}\Abracket{\pD_{\T^2}\varphi,\varphi}\dv_{\T^2}
    \coloneqq & 
    \Abracket{|\pD_{\T^2}|P^+\varphi - |\pD_{\T^2}|P^-\varphi ,\varphi}_{H^{-1/2}\times H^{1/2}} \\
    =&\| |\pD_{\T^2}|^{\frac{1}{2}}P^+\varphi\|_{L^2}^2 
     - \| |\pD_{\T^2}|^{\frac{1}{2}}P^-\varphi\|_{L^2}^2
\end{align}
where the last expression is well-defined variationally.  
In the same way, we have 
\begin{align}
    \tnorm{\varphi}_{H^{\frac{1}{2}}}^2 = \| |\pD_{\T^2}|^{\frac{1}{2}}\varphi \|_{L^2}^2 
    =&\| |\pD_{\T^2}|^{\frac{1}{2}}P^+\varphi\|_{L^2}^2 
     + \| |\pD_{\T^2}|^{\frac{1}{2}}P^-\varphi\|_{L^2}^2 \\
     =& \Abracket{|\pD_{\T^2}|P^+\varphi +|\pD_{\T^2}|P^-\varphi ,\varphi}_{H^{-1/2}\times H^{1/2}} \\
     \equiv & \int_{\T^2}\Abracket{|\pD_{\T^2}|\varphi,\varphi}\dv_{\T^2}
\end{align}
where the last integral is again understood in the sense of the duality pairing. 


\ 

Equivalently, since~$\Spect(\pD_{\T^2})$ has only point spectrum consisting of eigenvalues of finite multiplicity, one could use the Fourier expansions via eigenspinors and define~$|\pD_{\T^2}|^s$ and its domain as follows. 
List the eigenvalues of~$\pD_{\T^2}$ in an increasing order as 
\begin{align}
    -\infty \leftarrow \cdots \leq 
    \mu_{-k}\leq \cdots \leq \mu_{-1} <0 < \mu_1 \leq \cdots \leq \mu_k \leq \cdots \to +\infty
\end{align}
and let~$\braces{ \chi_k\mid k \in\mathbb{Z}_*}$ (where~$\mathbb{Z}_*$ denotes the set of nonzero integers) be a~$L^2$-orthonormal basis consisting of eigenspinors. 
Then any~$L^2$-section~$\varphi$ of the bundle~$\Sigma\T^2 \oplus \Sigma\T^2$ can be expanded in Fourier modes as 
\begin{align}
    \varphi=\sum_{k\in \mathbb{Z}_*} a_k \chi_k. 
\end{align}
Then~$\pD_{\T^2}$ acts on~$\varphi$ as
\begin{align}
    \pD_{\T^2}\varphi=\sum_{k\in\mathbb{Z}_*} \mu_k a_k \chi_k.
\end{align}
For any~$s\in\R$, we define~$H^s(\T^2,\C^4)$ by the Fourier expansion
\begin{align}
    H^s(\T^2,\C^4)
    =
    \braces{
        \varphi=\sum_{k\in\mathbb Z_*}a_k\chi_k
        \;\middle|\;
        \sum_{k\in\mathbb Z_*}|\mu_k|^{2s}|a_k|^2<+\infty
    },
\end{align}
where for~$s<0$ the space is understood as the completion of finite
Fourier sums with respect to the corresponding norm.
We define
\begin{align}
    |\pD_{\T^2}|^s\varphi
    =
    \sum_{k\in\mathbb Z_*}|\mu_k|^s a_k\chi_k,
\end{align}
with domain exactly~$H^{s}(\T^2,\C^4)$. Then
\begin{align}
    \tnorm{\varphi}_{H^s}
    =
    \||\pD_{\T^2}|^s\varphi\|_{L^2}
    =
    \parenthesis{
        \sum_{k\in\mathbb Z_*}|\mu_k|^{2s}|a_k|^2
    }^{\frac12},
\end{align}
for every~$s\in\R$.
In this way we can get a clearer picture of the fractional Sobolev spaces on~$\T^2$.
For~$s=\frac{1}{2}$, the subspace~$H^{\frac{1}{2},+}$ then consists of 
\begin{align}
    \varphi=\sum_{k>0} a_k \chi_k, & & \mbox{ with } & & \sum_{k>0}\mu_k |a_k|^2<+\infty,
\end{align}
and an analogous description for~$H^{\frac{1}{2},-}$ holds. 

 

\begin{rmk}\label{rmk:resolvent-estimate}
    We can now give another explanation of Lemma~\ref{lemma:invertibility-orthogonal}. 
    Using the Sobolev embeddings
    \begin{align}
    H^{\frac{1}{2}}(\T^2)\hookrightarrow L^4(\T^2,\C^4), & & L^{\frac{4}{3}}(\T^2,\C^4)\hookrightarrow H^{-\frac{1}{2}}(\T^2,\C^4),
    \end{align}
    it suffices to find a constant~$C>0$ such that
    \begin{align}
    \|\proj_{m,\zeta}^\perp \varphi\|_{H^{\frac{1}{2}}}\leq C\|\proj_{m,\zeta}^\perp(\pD_{\T^2}-a)\varphi\|_{H^{-\frac{1}{2}}}.
    \end{align}
    Recall that~$\proj_{m,\zeta}^\perp$ commutes with~$(\pD_{\T^2}-a)$.
    We thus restrict to spinors in the orthogonal complement of~$\Eigen(\pD_{\T^2};\sqrt{\zeta^2+m^2})$ and let 
    \begin{align}
    \varphi=\sum_{k\colon \mu_k\neq\sqrt{\zeta^2+m^2}}a_k \chi_k.
    \end{align}
    Then we have 
    \begin{align}
        \tnorm{\proj_{m,\zeta}^\perp(\pD_{\T^2}-a)\varphi }_{H^{-\frac{1}{2}}}^2
        =&\tnorm{ \sum_{k\colon \mu_k\neq\sqrt{\zeta^2+m^2}} a_k(\mu_k-a) \chi_k }_{H^{-\frac{1}{2}}}^2 \\
        =&\sum_{k\colon \mu_k\neq\sqrt{\zeta^2+m^2}}|\mu_k|^{-1}|(\mu_k-a)a_k|^2 \\
        =&\sum_{k\colon\mu_k\neq\sqrt{\zeta^2+m^2}} \left|\frac{\mu_k-a}{\mu_k}\right|^2 |\mu_k| |a_k|^2 \\
        \geq& \parenthesis{ \inf_{k\colon \mu_k\neq\sqrt{\zeta^2+m^2}} \left|\frac{\mu_k-a}{\mu_k}\right|^2  } \sum_{k\colon\mu_k\neq \sqrt{\zeta^2+m^2}} |\mu_k| |a_k|^2\\
        \geq& c_0^2 \tnorm{\proj_{m,\zeta}^\perp \varphi }_{H^{\frac{1}{2}}}^2.
    \end{align}
    In the last inequality, 
    \begin{align}
        0<c_0=1-\frac{\sqrt{\zeta^2+m^2}}{\sqrt{\zeta^2+m^2+\lambda_1^2}} 
        \leq 1-\frac{a}{\sqrt{\zeta^2+m^2+\lambda_1^2}}
        \leq \inf_{k\colon \mu_k\neq\sqrt{\zeta^2+m^2}} \left|\frac{\mu_k-a}{\mu_k}\right| 
    \end{align}
    for all~$a\in [0,\sqrt{\zeta^2+m^2}]$, where~$\lambda_1$ is the smallest positive eigenvalue of~$\pd_{\T^2}$ in~\eqref{eq:spectrum-pd}.
    This makes clear that the constant in Lemma~\ref{lemma:invertibility-orthogonal} is essentially~$c_0^{-1}$, multiplied by embedding constants.
\end{rmk}


\

With~$\zeta\in\bZ\frac{2\pi}{\ell_3}$ fixed, we are ready to consider the functional~$J_{\zeta}\colon H^{\frac{1}{2}}(\mathbb{T}^2;\C^4)\to \R$, defined by 
\begin{align}\label{eq:functional-T2}
    J_\zeta (\varphi)
    =&\int_{\T^2} \frac{1}{2}\Abracket{\pD_{\T^2} \varphi,\varphi }
     -\frac{a}{2}|\varphi|^2 -F(\varphi)\dv_{\T^2}
    \\
    =&\frac{1}{2}\parenthesis{ \tnorm{P^+\varphi}_{H^{1/2}}^2 
     - \tnorm{P^-\varphi }_{H^{1/2}}^2 
     -a\|\varphi\|_{L^2}^2 }
     -\int_{\T^2} F(\varphi)\dv_{\T^2}.
\end{align}
The Euler--Lagrange equation of~$J_\zeta$ is precisely given by~\eqref{eq:NDE-T2}.
We therefore seek nonzero critical points of~$J_{\zeta}$. 
However, the functional~$J_\zeta$ is strongly indefinite, and under the structural assumptions on the Soler-type nonlinearity considered here, the compactness required by the standard Palais--Smale or Cerami framework is not directly available. 
We therefore follow the perturbation strategy used in~\cite{WZ2026Stationary}: we first add a higher-order perturbation term, obtain critical points for the perturbed functionals, derive estimates that are uniform with respect to the perturbation parameter, and finally pass to the zero-perturbation limit.
Moreover, since~$\dim (\mathbb{T}^2)=2$ and~$H^{\frac{1}{2}}(\T^2)\hookrightarrow L^4(\T^2)$, we expect to deal with nonlinearities of subquartic growth. 



More precisely, we consider 
\begin{align}\label{eq:perturbed-functional}
    J_{\zeta,\eps}(\varphi)
    =J_\zeta(\varphi)- \eps\int_{\T^2} |\varphi|^{\alpha_2}\dv_{\T^2},
\end{align}
where~$\eps\in[0,1]$ is the perturbation parameter. 
The Euler--Lagrange equation for~$J_{\zeta,\eps}$ is 
\begin{align}\label{eq:perturbed-NDE}
    \pD_{\T^2}\varphi-a\varphi - \p F(\varphi)= \eps \alpha_2 |\varphi|^{\alpha_2 -2 }\varphi, \quad \mbox{ on } \T^2. 
\end{align}
It reduces to~\eqref{eq:NDE-T2} when~$\eps= 0$. 
For later convenience we abbreviate~$F(\varphi)+\eps|\varphi|^{\alpha_2}$ as~$F_\eps(\varphi)$, so that~\eqref{eq:perturbed-NDE} can be rewritten as 
\begin{align}
    \pD_{\T^2}\varphi = a\varphi + \p F_\eps(\varphi).
\end{align}
By enlarging the constants~$A_j$ if necessary, we may assume that the perturbed nonlinearities~$F_\eps$ satisfy (F1-4) uniformly, for any~$\eps\in [0,1]$. 
Since~$\alpha_2\geq\alpha_1>2$, they also satisfy 
\begin{align}
    \p_\varphi F_\eps(x,\varphi)[\varphi] \geq \alpha_1 F_\eps(x,\varphi).
\end{align}

We make a notational remark here: formally we should write the functional as~$J_{\zeta,\eps,a}$ since at the end of Section~\ref{sect:perturbed solutions} we will vary~$a$ in a suitable range. 
But before that, the temporal frequency is fixed.
Thus we decide to use the notation~$J_{\zeta,\eps}$ and~$J_{\zeta,\eps, a}$ interchangeably, with the latter preferred only when we need to emphasize the dependence on~$a$ and with the former adopted for simplicity of notation. 
This should cause no confusion. 


\section{Minimax solutions to the perturbed equations}\label{sect:perturbed solutions}

In this section we consider the perturbed functionals~$J_{\zeta,\eps}$ with~$\zeta\in\bZ\frac{2\pi}{\ell_3}$ and~$\eps\in (0,1]$, and obtain nonzero critical points of~$J_{\zeta,\eps}$ via a minimax method. 


We start with the verification of Palais--Smale condition for the perturbed functionals.
The argument goes in the same way as in~\cite{WZ2026Stationary} and we include it for completeness. 
Moreover, we observe that the possible critical levels of~$J_{\zeta,\eps}$ with~$\eps>0$ are necessarily nonnegative. 

\begin{prop}\label{prop:PS}
    Fix~$\eps\in (0,1]$ and let~$(\varphi_n)$ be a~$(PS)_c$ sequence for~$J_{\zeta,\eps}$.
    \begin{itemize}
        \item[(i)] The sequence~$(\varphi_n)$ is bounded in~$H^{\frac{1}{2}}$.
        \item[(ii)] There is a convergent subsequence of~$(\varphi_n)$ whose limit is a solution of~\eqref{eq:perturbed-NDE}.
        \item[(iii)] $c\geq 0$. Moreover, if~$c=0$, then the limit in (ii) is necessarily trivial.
    \end{itemize}
\end{prop}

\begin{proof}
    That~$(\varphi_n)$ is a~$(PS)_c$ sequence means
    \begin{align}\label{eq:PS level}
        J_{\zeta,\eps}(\varphi_n)
        =J_\zeta(\varphi_n)- \eps\int_{\T^2} |\varphi_n|^{\alpha_2}\dv_{\T^2} \to c,
    \end{align}
    and
    \begin{align}\label{eq:PS differential}
        \dd J_{\zeta,\eps}(\varphi_n)=\pD_{\T^2}\varphi_n-a\varphi_n-\partial F_\eps(\varphi_n)\; \to \; 0
         \quad \mbox{ in } H^{-\frac{1}{2}},
    \end{align}
    as~$n\to+\infty$.


    \

    (i) Recall that~$\partial F_\eps(\varphi_n): = \partial F(\varphi_n) + \eps\alpha_2|\varphi_n|^{\alpha_2-2}\varphi_n$. 
    We test~\eqref{eq:PS differential} against~$\varphi_n$ and then subtract it from the twice of~\eqref{eq:PS level}:
    \begin{align}\label{eq: the energy functional minus the F-derivative}
    2J_{\zeta,\eps}(\varphi_n)-\dd J_{\zeta,\eps}(\varphi_n)[\varphi_n]
    =&\int_{\T^2}  \partial F(\varphi_n)[\varphi_n]-2F(\varphi_n)+\eps (\alpha_2-2)|\varphi_n|^{\alpha_2}\dv_{\T^2} \\
    \geq & \eps(\alpha_2-2)\int_{\T^2}  |\varphi_n|^{\alpha_2}\dv_{\T^2},
    \end{align}
    where (F3) is used in the last inequality. Thus
    \begin{align}
    \int_{\T^2}  |\varphi_n|^{\alpha_2}\dv_{\T^2} \leq
    \frac{1}{\eps(\alpha_2-2)} \parenthesis{ 2c + o(1) + o(\|\varphi_n\|_{H^{\frac12}})}.
    \end{align}
    Observe that in the equation
    \begin{align}
    \pD_{\T^2} \varphi_n = a\varphi_n + \partial F(\varphi_n) + \eps \alpha_2 |\varphi_n|^{\alpha_2-2}\varphi_n +\dd J_{\zeta,\eps}(\varphi_n),
    \end{align}
    the right-hand side can be controlled by
    \begin{align}
    \| RHS\|_{H^{-\frac{1}{2}}}
    \leq & a \|\varphi_n\|_{H^{-\frac{1}{2}}} +C(A_2,\alpha_2)( 1+\||\varphi_n|^{\alpha_2-1}\|_{H^{-\frac{1}{2}}}) + o(1) \\
    \leq & C(a,\zeta, A_2, \eps, \alpha_2, \vol(\T^2))(1+ \|\varphi_n\|_{L^{\alpha_2}}^{\alpha_2-1})   + o(1) \\
    \leq&  C(a,\zeta, A_2, \eps, \alpha_2, \vol(\T^2)) \parenthesis{1+\frac{1}{\eps(\alpha_2-2)} \parenthesis{ 2c + o(1) + o(\|\varphi_n\|_{H^{\frac12}})}}^{\frac{\alpha_2-1}{\alpha_2}}.
    \end{align}
    Using the invertibility of~$\pD_{\T^2}$ we see that 
    \begin{align}
    \|\varphi_n\|_{H^{\frac{1}{2}}} \leq C \|\pD_{\T^2}\varphi_n\|_{H^{-\frac{1}{2}}}\leq C(c,a,\zeta, A_2, \eps,\alpha_2,\vol(\T^2)).
    \end{align}

    \

    (ii) By Banach--Alaoglu theorem, passing to a subsequence if necessary, we may assume that~$\varphi_n$ converges  to a limit~$\varphi_\eps\in H^{\frac{1}{2}}$ \emph{weakly} in~$H^{\frac{1}{2}}$ and strongly in~$L^q$ for any~$q<4$.
    In particular,~$\varphi_n$ converges strongly to~$\varphi_\eps$ in~$L^{\alpha_2}$ since~$\alpha_2 <4$. 
    We next show that~$\varphi_n\to\varphi_\eps$ strongly in~$H^{\frac{1}{2}}$ and~$\varphi_\eps$ is a critical point of~$J_{\zeta,\eps}$.

    The continuity assumption (F2) and growth assumption~\eqref{eq:F2-all} then imply 
    \begin{align}
        \p_\varphi F(\cdot, \varphi_n) \to \p_\varphi F(\cdot,\varphi_\eps), & & 
        |\varphi_n|^{\alpha_2-2}\varphi_n \to |\varphi_\eps|^{\alpha_2-2} \varphi_\eps, & & 
        \mbox{ in  } L^{\frac{\alpha_2}{\alpha_2-1}}.
    \end{align}
    Note that~$\frac{\alpha_2}{\alpha_2-1}>\frac{4}{3}$ and~$L^{\frac{4}{3}}(\T^2,\C^4) \hookrightarrow H^{-\frac{1}{2}}(\T^2,\C^4)$.
    Thus we may pass to the limit in~\eqref{eq:PS differential} and obtain that~$\dd J_{\zeta,\eps}(\varphi_\eps)=0$. 
    
    

    The difference~$\varphi_n-\varphi_\eps$ satisfies
    \begin{align}
        \pD_{\T^2}(\varphi_n-\varphi_\eps)
        =a(\varphi_n-\varphi_\eps)+\partial F(\varphi_n)-\partial F(\varphi_\eps) + \eps \alpha_2 \parenthesis{|\varphi_n|^{\alpha_2-2} \varphi_n - |\varphi_\eps|^{\alpha_2-2}\varphi_\eps }  + \dd J_{\zeta,\eps}(\varphi_n)
    \end{align}
    and the right-hand side above converges to~$0$ in~$H^{-\frac{1}{2}}$. 
    Again by the invertibility of~$\pD_{\T^2}$ we get 
    \begin{align}
        \|\varphi_n-\varphi_\eps\|_{H^{\frac{1}{2}}} \leq C\|\pD_{\T^2}(\varphi_n-\varphi_\eps)\|_{H^{-\frac{1}{2}}} \to 0.
    \end{align}
    Thus~$\varphi_n \to \varphi_\eps$ is actually strong in~$H^{\frac{1}{2}}$. 

    Therefore, the perturbed functional~$J_{\zeta,\eps}$ satisfies the Palais--Smale condition as long as~$\eps\in (0,1]$.

    \

    (iii) Now we have seen that the sequence~$(\varphi_n)$ is bounded in~$H^{\frac{1}{2}}$ and converges in~$H^{\frac{1}{2}}$ to~$\varphi_\eps$. 
    Thus
    \begin{align}
        c=&\lim_{n\to+\infty} J_{\zeta,\eps}(\varphi_n) - \frac{1}{2}\dd J_{\zeta,\eps}(\varphi_n)[\varphi_n] \\
        =&\lim_{n\to+\infty} \int_{\T^2} -F(\varphi_n)+\frac{1}{2}\partial F(\varphi_n)[\varphi_n]\dv_{\T^2}
           +\eps\int_{\T^2} -|\varphi_n|^{\alpha_2} +\frac{1}{2}\alpha_2|\varphi_n|^{\alpha_2}\dv_{\T^2} \\
        \geq & \lim_{n\to+\infty}\frac{1}{2}\eps(\alpha_2-2)\int_{\T^2} |\varphi_n|^{\alpha_2}\dv_{\T^2} \geq 0,
    \end{align}
    because of (F3) and~$\alpha_2>2$.
    Hence all possible critical levels of~$J_{\zeta,\eps}$ are nonnegative.
    Moreover, if~$c=0$, then~$\|\varphi_n\|_{L^{\alpha_2}}\to 0$ and the strong limit is~$\varphi_\eps$. 
    This holds for any $(PS)_0$ sequence, thus the only critical point at the level~$0$ is the zero spinor. 
\end{proof}


Next we search for minimax levels of the perturbed functional.
To this aim we need to find suitable linking structure of the functional~$J_{\zeta,\eps}$.
Although the functional is strongly indefinite, the explicit spectral description of~$\pD_{\T^2}$ provides us with local linking geometry. 


Consider the eigenspace of the lowest positive eigenvalue
\begin{align}
        \mathscr{P}= \Span_{\C}\braces{ \begin{pmatrix} 1 \\ 0 \\\frac{\zeta}{ m+\sqrt{\zeta^2+m^2} } \\ 0  \end{pmatrix},
        \begin{pmatrix} 0 \\ 1 \\ 0 \\ \frac{-\zeta}{m+\sqrt{\zeta^2+m^2}} \end{pmatrix} },
\end{align}
which is a vector space of complex dimension two.
It consists of parallel spinors on~$\T^2$, i.e. the component functions are constant.
But unlike the corresponding spinors in~\cite{WZ2026Stationary}, these vectors do not lie in the (+1)-eigenspace of~$\gamma^0$, which is crucial in the construction of links there.
However, we can still see that the upper components occupy at least a fixed portion of the norms of spinors, which is enough for our application.

Given~$\zeta\in\bZ\frac{2\pi}{\ell_3}$, we will use the notation
\begin{align}
    \theta\equiv \theta(\zeta,m)=1-\frac{m^2}{36(\zeta^2+m^2)}\in [\frac{35}{36},1)
\end{align}
throughout this section. 

\begin{lemma}\label{lemma:local control}
    For any~$\mu \in  (\theta\sqrt{\zeta^2+m^2},\sqrt{\zeta^2+m^2})$, and for any~$k\in\braces{1,2,3,4}$, there exists a real~$k$-dimensional space~$\mathscr{E}_{k}\subset H^{\frac{1}{2},+}$ and a positive number~$R>0$ such that for any~$\varphi\in H^{\frac{1}{2}}$, and for any~$\eps\in [0,1]$,
    \begin{multline}
        \parenthesis{ \tnorm{P^-\varphi}_{H^{\frac{1}{2}}} \leq R, \; \mbox{ and }\;  P^+\varphi\in\mathscr{E}_{k}, \, \tnorm{P^+\varphi }_{H^{\frac{1}{2}}}=R }  \\
        \quad \Longrightarrow  \quad
        \int_{\T^2} \frac{1}{2}\Abracket{\pD_{\T^2}\varphi,\varphi}-F_\eps(\varphi) \dv_{\T^2} \leq \int_{\T^2} \frac{\mu}{2}|\varphi|^2 \dv_{\T^2}.
    \end{multline}
\end{lemma}

\begin{proof}
Let~$\mathscr{E}_k$ be a real~$k$-dimensional subspace of~$\mathscr{P}$ consisting of constant spinors.
Then for any~$e\in\mathscr{E}_k$,~$\pD_{\T^2} e=\sqrt{\zeta^2+ m^2} e$.
Moreover,
\begin{align}
    \bar{e}e = \Abracket{\gamma^0 e, e}
    =& |e_{\up}|^2 - |e_{\down}|^2
    = |e_{\up}|^2 - \parenthesis{ \frac{\zeta}{m+\sqrt{\zeta^2+m^2}} }^2 |e_{\up}|^2 \\
    =&\parenthesis{  1 - \frac{\zeta^2}{(m+\sqrt{\zeta^2+m^2})^2} } |e_{\up}|^2
    = \frac{2 m}{m+\sqrt{\zeta^2+m^2}} |e_{\up}|^2,
\end{align}
\begin{align}
    |e|^2=|e_{\up}|^2 +|e_{\down}|^2
    =\parenthesis{  1 + \frac{\zeta^2}{(m+\sqrt{\zeta^2+m^2})^2} } |e_{\up}|^2
    =\frac{2\sqrt{\zeta^2+m^2}}{m+\sqrt{\zeta^2+m^2}}|e_{\up}|^2.
\end{align}
Thus
\begin{align}
    \bar{e}e = \frac{m}{\sqrt{\zeta^2 + m^2}} |e|^2.
\end{align}

Then for~$e=P^+\varphi\in\mathscr{E}_k$ with~$\tnorm{P^+\varphi}_{H^{1/2}} = \tnorm{e}_{H^{1/2}}=R$, we have
\begin{align}
    R^2
    =\int_{\T^2} \Abracket{P^+e,\pD_{\T^2} P^+ e}\dv_{\T^2} = \sqrt{\zeta^2+m^2} \int_{\T^2} |e|^2\dv_{\T^2}
\end{align}
and in particular
\begin{align}
    \int_{\T^2} \frac{1}{2}\Abracket{e,\pD_{\T^2} e}\dv_{\T^2}
    = \int_{\T^2} \frac{\sqrt{\zeta^2+m^2} }{2} |e|^2 \dv_{\T^2}
    = \frac{R^2}{2},
\end{align}
\begin{align}
    \int_{\T^2} \bar{e}e\dv_{\T^2} = \frac{m}{\sqrt{\zeta^2 + m^2}} \int_{\T^2} |e|^2 \dv_{\T^2}
    = \frac{m}{{\zeta^2 + m^2}}\cdot R^2.
\end{align}
For a general~$\varphi=P^-\varphi+e\in H^{\frac{1}{2},-}\oplus\mathscr{E}_k$ with~$\tnorm{P^+\varphi}_{H^{1/2}}=\tnorm{e}_{H^{1/2}}=R$ and~$\tnorm{P^-\varphi}_{H^{1/2}}\leq R$, we consider the following two cases.
\begin{itemize}
    \item If~$-\int_{\T^2}\Abracket{P^-\varphi,\pD_{\T^2} P^-\varphi}\dv_{\T^2}\geq(1-\frac{\mu}{\sqrt{\zeta^2+m^2}})R^2$, then, since~$F_\eps(\varphi)\geq 0$, we have
    \begin{align}
        \int_{\T^2} \frac{1}{2}\Abracket{\pD_{\T^2}\varphi,\varphi}-F_\eps(\varphi) \dv_{\T^2}
        =& \int_{\T^2} \frac{1}{2}\Abracket{P^+\varphi,\pD_{\T^2} P^+\varphi} +\frac{1}{2} \Abracket{P^-\varphi,\pD_{\T^2} P^-\varphi} -F_\eps(\varphi) \dv_{\T^2} \\
         \leq&          \frac{1}{2}\||\pD_{\T^2}|^{\frac12}P^+\varphi\|^2_{L^2} - \frac{1}{2}\||\pD_{\T^2}|^{\frac12}P^-\varphi\|^2_{L^2} \\
        \leq & \frac{1}{2}\cdot R^2 - \frac{1}{2}\cdot(1-\frac{\mu}{\sqrt{\zeta^2+m^2}})R^2\\
        = & \frac{\mu}{2}\cdot\frac{R^2}{\sqrt{\zeta^2+m^2}}  = \frac{\mu}{2}\int_{\T^2} |P^+ \varphi|^2 \dv_{\T^2}  \\
        \leq& \frac{\mu}{2}\int_{\T^2} |\varphi|^2\dv_{\T^2}.
    \end{align}

    \item If, on the other hand, ~$-\int_{\T^2}\Abracket{P^-\varphi,\pD_{\T^2} P^-\varphi}\dv_{\T^2}\leq(1-\frac{\mu}{\sqrt{\zeta^2+m^2}})R^2$, then 
    \begin{align}
    \|P^-\varphi\|_{L^2}^2
    \leq& -\frac{1}{\sqrt{\zeta^2+m^2}} \int_{\T^2} \Abracket{P^-\varphi,\pD_{\T^2} P^-\varphi }\dv_{\T^2} \\
    \leq& \frac{1}{\sqrt{\zeta^2+m^2}}(1-\frac{\mu}{\sqrt{\zeta^2+m^2}})R^2.
    \end{align}
    The negative quadratic contribution alone is insufficient.
    We turn to the nonlinearity to control the positive contribution.
    Indeed, using (F4), we see that
    \begin{align}
    -\int_{\T^2} F_\eps(\varphi)\dv_{\T^2}
    \leq& -\int_{\T^2} A_3 |\bar{\varphi}\varphi|^\nu -A_4 \dv_{\T^2}
    =-A_3\int_{\T^2} |\bar{\varphi}\varphi|^\nu\dv_{\T^2} + A_4 \cdot\vol(\T^2).
    \end{align}
    The term involving the Lorentz scalar~$\bar{\varphi}\varphi$ can be estimated by
    \begin{align}
    &\parenthesis{\int_{\T^2} |\bar{\varphi}\varphi|^\nu\dv_{\T^2} }^{\frac{1}{\nu}}
    \parenthesis{\int_{\T^2} 1 \dv_{\T^2}}^{1-\frac{1}{\nu}}\\
    \geq&  \int_{\T^2} \bar{\varphi}\varphi\dv_{\T^2}
    = \int_{\T^2} \Abracket{\gamma^0\parenthesis{P^+\varphi+ P^-\varphi}, P^+\varphi + P^-\varphi}\dv_{\T^2} \\
    =& \int_{\T^2} \Abracket{\gamma^0 P^+ \varphi, P^+\varphi} + 2\operatorname{Re}\Abracket{ e,\gamma^0P^-\varphi}
    + \Abracket{\gamma^0 P^- \varphi, P^-\varphi}\dv_{\T^2} \\
    \geq & \frac{mR^2}{{\zeta^2+m^2}} -2\|e\|_{L^2}\|P^-\varphi\|_{L^2}-\|P^-\varphi\|_{L^2}^2\\
    \geq&\frac{mR^2}{{\zeta^2+m^2}}-2\parenthesis{\frac{R^2}{\sqrt{\zeta^2+m^2}} }^{\frac{1}{2}}\parenthesis{ \parenthesis{1-\frac{\mu}{\sqrt{\zeta^2+m^2}}}\frac{R^2}{\sqrt{\zeta^2+m^2}}  }^{\frac{1}{2}}  \\
        & -\parenthesis{1-\frac{\mu}{\sqrt{\zeta^2+m^2}}}\frac{R^2}{\sqrt{\zeta^2+m^2}} \\
    =&\frac{R^2}{{\zeta^2+m^2}} \parenthesis{m+\mu-\sqrt{\zeta^2+m^2} -2(\zeta^2+m^2)^{\frac{1}{4}}\sqrt{\sqrt{ \zeta^2+m^2 }  -\mu} }.
\end{align}
To get a positive lower bound of the last line above, we let~$\mu=(1-\delta)\sqrt{\zeta^2+m^2}$ with~$\delta\in(0,1)$ and compute
\begin{align}
   &m+\mu-\sqrt{\zeta^2+m^2} -2(\zeta^2+m^2)^{\frac{1}{4}}\sqrt{\sqrt{ \zeta^2+m^2 }  -\mu}  \\
   =& m+(1-\delta)\sqrt{\zeta^2+m^2}-\sqrt{\zeta^2+m^2} -2(\zeta^2+m^2)^{\frac{1}{4}}\sqrt{\delta\sqrt{\zeta^2+m^2}}\\
   =& m -\delta\sqrt{ \zeta^2+m^2 }-2\sqrt{\delta}\sqrt{ \zeta^2+m^2 }\\
   \geq& m-3\sqrt{\delta}\sqrt{ \zeta^2+m^2 }\\
   \geq& \frac{m}{2} \qquad\qquad \mbox{ provided }\quad \delta\leq\frac{m^2}{36(\zeta^2+m^2)}.
\end{align}
For such~$\mu$, 
\begin{align}
    \int_{\T^2} |\bar{\varphi}\varphi|^\nu \dv_{\T^2}
    \geq  \frac{1}{\vol(\T^2)^{\nu -1}} \parenthesis{\frac{m}{2}\cdot\frac{R^2}{\zeta^2+m^2}}^{\nu}.
\end{align}
In this case we have
\begin{align}
    \int_{\T^2} & \frac{1}{2}\Abracket{\varphi,\pD_{\T^2} \varphi}- F_\eps(\varphi) \dv_{\T^2} \\
    =& \int_{\T^2} \frac{1}{2}\Abracket{P^+\varphi,\pD_{\T^2} P^+\varphi} \dv_{\T^2} + \int_{\T^2}\frac{1}{2}\Abracket{P^-\varphi,\pD_{\T^2}P^-\varphi} \dv_{\T^2} - \int_{\T^2} F_\eps(\varphi)\dv_{\T^2} \\
    =&\frac{1}{2}\parenthesis{ \tnorm{P^+\varphi }_{H^{\frac{1}{2}}}^2 - \tnorm{P^-\varphi}_{H^{\frac{1}{2}}}^2  }
     -\int_{\T^2} F_\eps(\varphi)\dv_{\T^2}\\
    \leq & \frac{1}{2}R^2
    - \frac{A_3m^\nu R^{2\nu}}{\vol(\T^2)^{\nu-1} 2^\nu({\zeta^2+m^2})^\nu}
     +  A_4 \vol(\T^2).
\end{align}
Since~$\nu>1$ and~$A_3>0$, we can choose~$R=R(\vol(\T^2), \zeta,m, A_3, A_4, \nu)\gg 1$ such that the right-hand side above is nonpositive.
For example, it suffices to have 
\begin{align}
    R=\max\braces{\sqrt{2A_4 \vol(\T^2)} , \; \parenthesis{\frac{2^\nu}{A_3}}^{\frac{1}{2(\nu-1)}} \sqrt{\vol(\T^2)} \parenthesis{\frac{\zeta^2+m^2}{m}}^{\frac{\nu}{2(\nu-1)}} }. 
\end{align}
\end{itemize}
The desired conclusion follows.
\end{proof}
Note that the choices of~$\mathscr{E}_k$ and~$R$ are independent of~$\eps\in [0,1]$ and also independent of the frequency~$a\in \R$, a property crucial for later analysis.

\begin{rmk}
    It is possible to get a sharper estimate for the admissible values of~$\mu$.
    Indeed, consider the function~$\eta\colon [0,\sqrt{\zeta^2+m^2}]\to\R$ given by
    \begin{align}
    \eta(\mu)\coloneqq m+\mu-\sqrt{\zeta^2+m^2} -2(\zeta^2+m^2)^{\frac{1}{4}}\sqrt{\sqrt{ \zeta^2+m^2 }  -\mu}.
    \end{align}
    Then~$\eta$ is continuous, increasing and~$\eta(\sqrt{\zeta^2+m^2})=m>0$,~$\eta(\sqrt{\zeta^2+m^2}-m)<0$.
    Thus there exists~$\mu_0> \sqrt{\zeta^2+m^2}-m $ such that~$\eta|_{[\mu_0,\sqrt{\zeta^2+m^2}]}>\frac{m}{2}$.
    Furthermore, the unique root of~$\eta$ is actually given by
    \begin{align}
        \mu_*=\sqrt{\zeta^2+m^2}-\parenthesis{\frac{m}{\sqrt{\sqrt{\zeta^2+m^2}+m}+ (\zeta^2+m^2)^{1/4} } }^2.
    \end{align}
    We take the clean and explicit value of~$\theta$ above because it provides a convenient uniform interval for the linking argument.
    For our purpose in the variational analysis, it suffices to have a nonempty range of~$\mu$ where the conclusion of Lemma~\ref{lemma:local control} holds. 
    

    
\end{rmk}


For~$k=1$ we take $\mathscr{E}_1=\Span(e)$ with~$\tnorm{e}_{H^{\frac{1}{2}}}=1$ and let
\begin{align}
    \mathscr{C}_1(R)
    \coloneqq \braces{\varphi= P^-\varphi + s e\in H^{1/2}(\T^2,\C^4) \mid  \tnorm{P^-\varphi}_{H^{1/2}}\leq R, \; s \in [0, R]}.
\end{align}
To get a linking structure we consider the candidate
\begin{align}
    \mathscr{S}^+(r)\coloneqq \braces{\varphi\in H^{\frac{1}{2},+} \; \mid \; \tnorm{\varphi}_{H^{1/2}}=r }.
\end{align}
for some~$r\in(0,R)$ to be determined later.
Similar to the situation of~\cite{WZ2026Stationary}, we have
\begin{lemma}\label{lemma:negativity on boundary-1}
    Suppose~$a \in (\theta\sqrt{\zeta^2+m^2},\sqrt{\zeta^2+m^2})$. 
    Then~$J_{\zeta,\eps}|_{\p\mathscr{C}_1(R)}\leq 0$ for any~$\eps\in [0,1]$.
\end{lemma}
\begin{lemma}\label{lemma:positivity in interior-1}
    There exists~$r\in (0,R)$ and~$C_*>0$ such that for any~$\eps\in [0,1]$, we have
    \begin{align}
         J_{\zeta,\eps}|_{\mathscr{S}^+(r)}\geq C_*>0.
    \end{align}
\end{lemma}


\begin{proof}[Proof of Lemma~\ref{lemma:negativity on boundary-1}]
    Let~$\varphi\in\p\mathscr{C}_1(R)$, namely, either~$\tnorm{P^-\varphi}_{H^{1/2}}=R$ or~$s\in \braces{0,R}$.
    \begin{itemize}
        \item If~$\tnorm{P^-\varphi}_{H^{1/2}}=R$, then for any~$s\in[0,R]$, using~$F_\eps(\varphi)\geq 0$,
            \begin{align}
                J_{\zeta,\eps}(P^-\varphi+s e)&\leq \frac{1}{2}(1-\frac{a}{\sqrt{\zeta^2+m^2}})s^2 -  \frac{1}{2}\tnorm{P^-\varphi}^2_{H^{1/2}} \\
                &\leq \frac{1}{2}(1-\frac{a}{\sqrt{\zeta^2+m^2}})R^2 -\frac{1}{2}R^2< 0.
            \end{align}
        \item If~$s=0$, then~$J_{\zeta,\eps}(P^-\varphi)\leq 0$ since~$F_\eps(\varphi)\geq 0$; and if~$s=R$, then~$\tnorm{P^-\varphi}_{H^{1/2}}\leq R$ and~$\tnorm{P^+\varphi}_{H^{1/2}}= R\tnorm{e}_{H^{1/2}}=R$, applying Lemma~\ref{lemma:local control} with~$\mu \in (\theta\sqrt{\zeta^2+m^2}, a]$, we again have~$J_{\zeta,\eps}(\varphi)\leq 0$.
    \end{itemize}
\end{proof}

\begin{proof}[Proof of Lemma~\ref{lemma:positivity in interior-1}]
Note that for~$\varphi=P^+\varphi\in H^{\frac{1}{2},+}$,
\begin{align}
    \|\varphi\|^2_{L^2}\leq \frac{1}{\sqrt{m^2+\zeta^2}}\tnorm{\varphi}^2_{H^{\frac{1}{2}}}.
\end{align}
By (F1) we have 
\begin{align}
    J_{\zeta,\eps}(\varphi)
    \geq &  \frac{1}{2}\int_{\mathbb{T}^2} \Abracket{\varphi, \pD_{\T^2}\varphi} - a|\varphi|^2 \dv_{\T^2}
        - A_1\int_{\mathbb{T}^2} |\varphi|^{\alpha_1} + |\varphi|^{\alpha_2}\dv_{\T^2}\\
    \geq& \frac{1}{2}(1-\frac{a}{\sqrt{m^2+\zeta^2}})\tnorm{\varphi}^2_{H^{\frac{1}{2}}}
    -C(\tnorm{\varphi}^{\alpha_1}_{H^{\frac{1}{2}}}+\tnorm{\varphi}^{\alpha_2}_{H^{\frac{1}{2}}}).
\end{align}
Hence for sufficiently small~$r=\tnorm{\varphi}_{H^{1/2}}>0$, depending only on the fixed data~$m,\zeta,A_1, \vol(\T^2),\alpha_1$ and on~$a$,
\begin{align}
    \inf_{\mathscr{S}^+(r)} J_{\zeta,\eps}(\varphi^+)>0.
\end{align}
\end{proof}

We note that the~$r$ and~$C_*$ can be chosen to be independent of~$\eps\in [0,1]$.
Furthermore,
\begin{align}\label{eq:intersection}
    \mathscr{C}_1(R)\cap \mathscr{S}^+(r)= \braces{ r e} \neq\emptyset.
\end{align}
It follows that
\begin{align}\label{eq:level estimate}
    \sup_{\p \mathscr{C}_1(R)} J_{\zeta,\eps} = 0 
    < C_* 
    \leq  \inf_{\mathscr{S}^+(r)} J_{\zeta,\eps} 
    \leq \sup_{\mathscr{C}_1(R)} J_{\zeta,\eps} <+\infty
\end{align}
for each~$\eps\in [0,1]$, where the last upper bound holds due to the Sobolev embedding together with the fact~$\alpha_2<4$.
That is,~$\p\mathscr{C}_1(R)$ and~$\mathscr{S}^+(r)$ provide the local linking geometry for~$J_{\zeta,\eps}$. 

\

Next we consider deformations of the set~$\mathscr{C}_1(R)$ and try to define the minimax levels in a suitable way. 
Since~$F$ is only assumed to be~$C^{1,\alpha}_{loc}$, the gradient of~$J_{\zeta,\eps}$ is not guaranteed to be locally Lipschitz, hence we cannot directly use the negative gradient flow of~$J_{\zeta,\eps}$. 
The same issue arises in~\cite{WZ2026Stationary} and is resolved there using \emph{admissible pseudo-gradient fields and flows}. 
Those can also be used here on the two-dimensional torus. 
We briefly recall (and adapt suitably) the main constructions before we define the minimax levels. 


A vector field~$W$ on~$H^{\frac{1}{2}}(\mathbb{T}^2,\C^4)$ is called an \emph{admissible pseudo-gradient field} if it has the form 
\begin{align}
    W(\varphi)= \eta(\varphi) \parenthesis{ (P^+ -P^- )\varphi - \widetilde{K}(\varphi)},
\end{align}
where
\begin{itemize}
    \item[(W1)] $\eta\colon H^{\frac{1}{2}}\to [0,1]$ is a locally Lipschitz function,
    \item[(W2)] $\widetilde{K}\colon H^{\frac{1}{2}}\to H^{\frac{1}{2}}$ is locally Lipschitz and is compact on bounded sets;
    \item[(W3)] $W|_{\p\mathscr{C}_1(R)} = 0$;
    \item[(W4)] On the region~$\braces{\varphi\colon \eta(\varphi)>0}$, there holds~$\dd J_{\zeta,\eps}(\varphi)[P^+\varphi -P^- \varphi -\widetilde{K}(\varphi)]\geq 0$;
    \item[(W5)] $\tnorm{W(\varphi)}_{H^{1/2}} \leq 1$ for any~$\varphi\in H^{\frac{1}{2}}$. 
\end{itemize}
To such a vector field we associate an admissible pseudo-gradient flow~$\phi^W_t\colon H^{\frac{1}{2}}\to H^{\frac{1}{2}}$ generated by~$-W$: 
\begin{align}
    \begin{cases}
        \frac{\p\phi^W_t(\varphi)}{\p t} = - W(\phi^W_t(\varphi)), & \quad \forall \varphi \in H^{\frac{1}{2}}, \; \forall t\geq 0, \\
        \phi^W_0(\varphi)=\varphi, & \quad \forall \varphi\in H^{\frac{1}{2}}.
    \end{cases}
\end{align}
The flow fixes the boundary~$\p\mathscr{C}_1(R)$ pointwise and does not increase the functional~$J_{\zeta,\eps}$.  
For each~$t\geq 0$, we call~$\phi^W_t$ a time-$t$ map of the admissible pseudo-gradient flow. 
There are plenty of such admissible pseudo-gradient vector fields. 
Moreover, using variation of constants, we see that along the flow the negative part has the form
\begin{align}
    P^-\phi^W_t(\varphi)
    = e^{u^W(t,\varphi)}P^-\varphi + C^W(t,\varphi)
\end{align}
with 
\begin{align}
    u^W(t,\varphi)=\int_0^t \eta(\phi^W_s(\varphi))\dd s, \qquad  0\leq u^W(t,\varphi)\leq t ,
\end{align}
\begin{align}
    C^W(t,\varphi)
    = \int_0^t e^{\int_s^t \eta(\phi^W_\tau(\varphi))\dd \tau } \eta(\phi^W_s(\varphi))P^-\widetilde{K}(\phi^W_s(\varphi))\dd{s}. 
\end{align}
In particular, for any bounded set~$B\subset H^{\frac{1}{2}}$ and any bounded time interval~$[0,T]$, the map 
\begin{align}
    C^W\colon [0,T]\times B \to H^{\frac{1}{2},-}
\end{align}
is continuous and compact.


The deformation class is the collection of finite compositions of time maps of admissible pseudo-gradient flows:
\begin{align}
    \Gamma\coloneqq \braces{ \phi^{W_k}_{t_k}\circ \cdots\circ \phi^{W_1}_{t_1} \mid W_j \mbox{ is admissible }, t_j\geq 0, \; k\in\mathbb{N} }. 
\end{align}
It is nonempty and closed under composition. 
Each element~$\tilde{\phi}$ is homotopic to the identity map by concatenating finitely many flow segments. 
We denote such a homotopy by~$\Phi(t)$:~$\Phi(0)=\id$,~$\Phi(1)=\tilde{\phi}$, and observe that for each~$t\in [0,1]$,
\begin{itemize}
    \item[(1)] $\Phi(t)$ fixes~$\p\mathscr{C}_1(R)$ pointwise and hence~$\Phi(t,\p\mathscr{C}_1(R))\cap \mathscr{S}^+(r) =\emptyset$;
    \item[(2)] $P^-\Phi(t,\varphi)=e^{u(t,\varphi)}P^-\varphi+ C_\Phi(t,\varphi)$ where~$u$ is continuous and bounded, and~$C_\Phi$ is continuous and compact;
    \item[(3)] $\Phi(t)(\mathscr{C}_1(R))$ intersects~$\mathscr{S}^+(r)$:~$\Phi(t,\mathscr{C}_1(R))\cap \mathscr{S}^+(r)\neq \emptyset$, see \cite[Lemma 4.6]{WZ2026Stationary}.
\end{itemize}


We can then define the following minimax level 
\begin{align}
    \Lambda_1(\zeta,\eps) \coloneqq \inf_{\tilde{\phi}\in \Gamma} \sup_{\tilde{\phi}(\mathscr{C}_1(R))} J_{\zeta,\eps}. 
\end{align}
By the standard minimax deformation argument associated with the admissible class~$\Gamma$ (see~\cite{WZ2026Stationary} for details), there exists a Palais--Smale sequence at the level~$\Lambda_1(\zeta,\eps)>0$. 
Indeed, otherwise an admissible pseudogradient deformation would lower the maximal value below~$\Lambda_1(\zeta,\eps)$ contradicting the definition of the minimax level. 
Since the perturbed functional~$J_{\zeta,\eps}$ satisfies the Palais--Smale condition by Proposition~\ref{prop:PS}, this sequence has a convergent subsequence whose limit is a critical point of~$J_{\zeta,\eps}$ at level~$\Lambda_1(\zeta,\eps)$, confirming that~$\Lambda_1(\zeta,\eps)$ is a critical level for~$J_{\zeta,\eps}$. 

We have seen that the critical levels satisfy
\begin{align}
    0<C_*\leq \Lambda_1(\zeta,\eps) \leq C^*<+\infty
\end{align}
with the lower and upper bounds independent of~$\eps\in [0,1]$. 
An explicit upper bound can be given by 
\begin{align}
    \sup_{\mathscr{C}_1(R)} J_{\zeta,\eps} 
    \leq \sup_{ \varphi\in \mathscr{C}_1(R)} \frac{1}{2}\tnorm{P^+\varphi}_{H^{\frac{1}{2}}}^2 - \frac{a}{2}\|P^+\varphi\|_{L^2}^2 = \frac{1}{2}\parenthesis{1-\frac{a}{\sqrt{\zeta^2+m^2}}}R^2. 
\end{align}
Thus we can take, for example,
\begin{align}
    C^*(a)= \frac{1}{2}\parenthesis{1-\frac{a}{\sqrt{\zeta^2+m^2}}}R^2. 
\end{align}
Recall that~$R$ is independent of~$a$ and~$\eps$.
Thus we have 
\begin{align}\label{eq:upper-level-estimate}
    \lim_{a\to \sqrt{\zeta^2+m^2}-0} C^*(a)=0.
\end{align}

Therefore, we obtain
\begin{thm}\label{thm:existence for perturbed NDE}
    For each~$\eps\in (0,1]$ and~$a\in (\theta\sqrt{\zeta^2+m^2},\sqrt{\zeta^2+m^2})$, there exists a nontrivial solution~$\varphi_{\eps,a}$ to~\eqref{eq:perturbed-NDE} such that
    \begin{align}
        0<C_* \leq J_{\zeta,\eps}(\varphi_{\eps,a})=\Lambda_1(\zeta,\eps)\leq C^*(a)<+\infty,
    \end{align}
\end{thm}

We next prove estimates uniform in~$\eps\in (0,1]$ after possibly restricting~$a$ to a smaller left neighborhood of~$\sqrt{\zeta^2+m^2}$. 
Since we aim to get an estimate which is also uniform in~$a$, for~$a$ in a suitable subset of~$(\theta\sqrt{\zeta^2+m^2},\sqrt{\zeta^2+m^2})$, from now on we denote the perturbed functional by~$J_{\zeta,\eps,a}$ to emphasize the dependence on~$a$. 


\begin{thm}\label{thm:uniform estimate}
    There exists~$m_*\in (\theta\sqrt{\zeta^2+m^2},\sqrt{\zeta^2+m^2})$ and~$C>0$, depending on the structure constants~$A_j$,~$\alpha_1,\alpha_2,\nu,\beta$, the physical parameters~$m,\zeta$ and the torus~$\T^2$, such that for any~$a\in (m_*, \sqrt{\zeta^2+m^2})$ and any~$\eps\in (0,1]$, 
    \begin{align}\label{eq:uniform estimate}
         0<\|\varphi_{\eps,a}\|_{H^1(\T^2,\C^4)}\leq C. 
    \end{align}
\end{thm}

\begin{proof}
    The proof follows the strategy of~\cite{WZ2026Stationary}, with suitable modifications reflecting the shifted lowest eigenspaces. 
    We will first bound the nonlinearity part in integral sense, then use a contradiction argument to show that the perturbed solutions have uniformly bounded~$L^4$ norms, and then use a bootstrap argument to get a uniform bound in~$H^1$. 

    \noindent\textbf{Step 1.} \emph{There exists a constant~$C=C(\alpha_1, C^*(a))>0$ such that
    \begin{align}
        0\leq \int_{\T^2} F_\eps(\varphi_{\eps,a})\, \dv_{\T^2}\leq C.
    \end{align}
    }


    We use~\eqref{eq:perturbed-NDE} and (F3) to get
    \begin{align}
        2J_{\zeta,\eps,a}(\varphi_{\eps,a}) + 2\int_{\T^2} F_\eps(\varphi_{\eps,a})\dv_{\T^2}
        =&\int_{\T^2} \Abracket{\varphi_{\eps,a},\pD_{\T^2}\varphi_{\eps,a}}-a|\varphi_{\eps,a}|^2\dv_{\T^2} \\
        =&\int_{\T^2} \partial F_\eps(\varphi_{\eps,a})[\varphi_{\eps,a}]\dv_{\T^2} \\
        \geq& \alpha_1\int_{\T^2} F_\eps(\varphi_{\eps,a}) \dv_{\T^2}.
    \end{align}
    As~$\alpha_1>2$ and~$J_{\zeta,\eps}(\varphi_{\eps,a})\in [C_*, C^*(a)]$, it follows that
    \begin{align}
        0\leq \int_{\T^2} F_\eps(\varphi_{\eps,a})\dv_{\T^2} \leq \frac{2C^*(a)}{\alpha_1-2}.
    \end{align}
    Since~$F_\eps$ is a sum of two nonnegative summands, we have 
    \begin{align}\label{eq:uniform estimate for perturbation}
    0\leq \int_{\T^2} F(\varphi_{\eps,a})\dv_{\T^2}\leq \frac{2C^*(a)}{\alpha_1-2}, & & 
    \mbox{ and } & & 
    0\leq \eps\int_{\T^2} |\varphi_{\eps,a}|^{\alpha_2}\dv_{\T^2}\leq \frac{2C^*(a)}{\alpha_1-2}.
    \end{align}
    Note that the above bound on the~$L^{\alpha_2}$ norms blows up as~$\eps\to0^+$.

    \


\noindent\textbf{Step 2.} \emph{We establish a uniform bound of the norms~$\|\varphi_\eps\|_{L^4}$.} 

Argue by contradiction, and suppose that there is no such uniform~$L^4$ bound on the perturbed solutions, no matter how close~$a$ is to the lowest positive eigenvalue~$\sqrt{\zeta^2+m^2}$. 
Then there would exist sequences~$(a_n)\subset (\theta\sqrt{\zeta^2+m^2},\sqrt{\zeta^2+m^2})$ and~$(\eps_n)\subset (0,1]$ such that
\begin{align}
    a_n\to\sqrt{\zeta^2+m^2}
    & & \mbox{ and } & &
    \lambda_n=\|\varphi_{\eps_n,a_n}\|_{L^4(\T^2)}\to+\infty. 
\end{align}
Passing to a subsequence if necessary, the sequence~$\eps_n$ also converges, say, to~$\eps_*\in [0,1]$. 
Note that 
\begin{align}\label{eq:NDE-an}
    (\pD_{\T^2} -a_n) \varphi_{\eps_n,a_n} 
    = \p F(\varphi_{\eps_n,a_n}) 
    + \alpha_2 \eps_n |\varphi_{\eps_n,a_n}|^{\alpha_2 -2} \varphi_{\eps_n,a_n}, 
\end{align}
\begin{align}
    0<J_{\zeta,\eps_n,a_n} (\varphi_{\eps_n,a_n})  \leq C^*(a_n) \to 0, \quad \mbox{ as } a_n \to \sqrt{\zeta^2+m^2}. 
\end{align}

\ 

\emph{Claim: $\eps_*=0$.}
\begin{proof}[Argument for Claim.] If~$\eps_*>0$, then we may assume that~$\eps_n>\frac{\eps_*}{2}>0$ for all~$n\geq 1$. 
    From~\eqref{eq:uniform estimate for perturbation} it follows that
    \begin{align}
        \|\varphi_{\eps_n,a_n}\|_{L^{\alpha_2}}^{\alpha_2} \leq \frac{4C^*(a_n)}{\eps_*(\alpha_1 -2)} \to 0.
    \end{align}
    By~\eqref{eq:NDE-an} and using~\eqref{eq:F2-all}, we get 
    \begin{align}
        \|\pD_{\T^2} \varphi_{\eps_n,a_n}\|_{L^{\frac{\alpha_2}{\alpha_2 -1}}(\T^2)}
        =& \| a_n \varphi_{\eps_n,a_n} +\p F(\varphi_{\eps_n,a_n}) 
    + \alpha_2 \eps_n |\varphi_{\eps_n,a_n}|^{\alpha_2 -2} \varphi_{\eps_n,a_n} \|_{L^{\frac{\alpha_2}{\alpha_2-1}}} \leq C<+\infty. 
    \end{align}
    Since~$\frac{\alpha_2}{\alpha_2 -1}>\frac{4}{3}$, and~$W^{1,\frac{4}{3}}(\T^2)\hookrightarrow L^4(\T^2)$, this gives a uniform bound of~$\|\varphi_{\eps_n,a_n}\|_{L^4}$, contradicting that~$\lambda_n\to+\infty$.
    Thus~$\eps_* =\lim\limits_{n\to+\infty} \eps_n =0$. 
\end{proof}


Normalize the solutions~$\varphi_{\eps_n,a_n}$ to
\begin{align}
    \vartheta_n\coloneqq \frac{\varphi_{\eps_n,a_n}}{\lambda_n}, & & \|\vartheta_n\|_{L^4}=1, & & n\geq 1,
\end{align}
which satisfy the equations
\begin{align}\label{eq:normalized eqn}
    \pD_{\T^2}\vartheta_n-a_n\vartheta_n=\frac{\p F(\varphi_{\eps_n,a_n})}{\lambda_n} + \alpha_2\eps_n |\varphi_{\eps_n,a_n}|^{\alpha_2 -2 }\vartheta_n, \qquad \mbox{ on } \; \mathbb{T}^2.
\end{align}
The perturbation terms converge to zero in~$L^{\frac{\alpha_2}{\alpha_2-1}}$:
\begin{align}
    \|\alpha_2\eps_n |\varphi_{\eps_n,a_n}|^{\alpha_2 -2 }\vartheta_n \|_{L^{\frac{\alpha_2}{\alpha_2 -1}}}
    = & \frac{\alpha_2 \eps_n^{\frac{1}{\alpha_2}} }{\lambda_n}
     \parenthesis{ \int_{\mathbb{T}^2} \left|\eps_n^{\frac{\alpha_2-1}{\alpha_2}} |\varphi_{\eps_n,a_n}|^{\alpha_2 -1} \right|^{\frac{\alpha_2}{\alpha_2-1}} \dv_{\T^2}   }^{\frac{\alpha_2-1}{\alpha_2}}\\
    \leq& \frac{\alpha_2 \eps_n^{\frac{1}{\alpha_2}} }{\lambda_n}
    \parenthesis{ \int_{\mathbb{T}^2} \eps_n |\varphi_{\eps_n,a_n}|^{\alpha_2}\dv_{\T^2}}^{\frac{\alpha_2-1}{\alpha_2}} \\
    \leq& \frac{\alpha_2 \eps_n^{\frac{1}{\alpha_2}} }{\lambda_n}\parenthesis{ \frac{2C^*(a_n)}{\alpha_1-2} }^{\frac{\alpha_2-1}{\alpha_2}} \to 0 ,\qquad \mbox{ as } n\to+\infty. 
\end{align}
For the nonlinear term involving~$\p F$, we need a more careful analysis.
The assumption (F5) implies the pointwise estimate (with~$\delta=1$)
\begin{align}
    \left|\frac{1}{\lambda_n} \p F(\varphi_{\eps_n,a_n})\right|
    \leq C \parenthesis{1+ F(\varphi_{\eps_n,a_n})^{\frac{1}{\beta}}} |\vartheta_n| \in L^{\frac{4\beta}{4+\beta}},
\end{align}
hence the norm estimate
\begin{align}
    \|\frac{1}{\lambda_n}\p F(\varphi_{\eps_n,a_n})\|_{L^{\frac{4\beta}{4+\beta}}}
    \leq C \parenthesis{ \vol(\mathbb{T}^2)+\frac{2 C^*(a_n)}{\alpha_1-2} }^{\frac{1}{\beta}}
\end{align}
due to \noindent\textbf{Step 1.}
Let 
\begin{align}
    p\coloneqq \min\braces{ \frac{4\beta}{4+\beta}, \frac{\alpha_2}{\alpha_2-1} } >\frac{4}{3}. 
\end{align}
Then~$ (\pD_{\T^2}-a_n)\vartheta_n \in L^p(\mathbb{T}^2)$ with
\begin{align}
    \|(\pD_{\T^2}-a_n)\vartheta_n\|_{L^p} \leq 2C \parenthesis{ \vol(\mathbb{T}^2)+\frac{2 C^*(a_n)}{\alpha_1-2} }^{\frac{1}{\beta}}.
\end{align}
It then follows that the sequence ~$(\vartheta_n)$ is uniformly bounded in~$W^{1,p}$ and,  upon passing to a subsequence, converges in~$L^4$ to some~$\vartheta_\infty$ with~$\|\vartheta_\infty\|_{L^4}=1$.
Moreover, by (F4) and~\eqref{eq:uniform estimate for perturbation}, 
\begin{align}
    \int_{\mathbb{T}^2} |\bar{\vartheta}_{n}\vartheta_{n}|^\nu\dv_{\T^2} \leq \frac{1}{A_3\lambda_n^{2\nu}} \parenthesis{ \int_{\mathbb{T}^2} F_{\eps_n}(\varphi_{\eps_n,a_n})\dv_{\T^2} + A_4\vol(\mathbb{T}^2)}
    \leq \frac{C+A_4 \vol(\mathbb{T}^2)}{A_3\lambda_n^{2\nu}}\to 0
\end{align}
as~$n\to+\infty$.
It follows that~$\bar{\vartheta}_\infty\vartheta_\infty\equiv 0$ almost everywhere. 
That is,~$\vartheta_\infty$ lies in the Lorentz null cone~$\mathcal{N}$,
so it cannot be in~$\Eigen(\pD_{\T^2};\sqrt{\zeta^2+m^2})$.
% Recall that 
% \begin{align}
%         \Eigen(\pD_{\T^2};\sqrt{\zeta^2+m^2})= \Span_{\C}\braces{ \begin{pmatrix} 1 \\ 0 \\\frac{\zeta}{ m+\sqrt{\zeta^2+m^2} } \\ 0  \end{pmatrix},
%         \begin{pmatrix} 0 \\ 1 \\ 0 \\ \frac{-\zeta}{m+\sqrt{\zeta^2+m^2}} \end{pmatrix} }.
% \end{align}
% This space consists of constant spinors on $\mathbb{T}^2$. 
% Such spinors no longer lie in the $(+1)$-eigenspace of $\gamma^0$, a property used in the contradiction argument in~\cite{WZ2026Stationary}.
% Nevertheless, they cannot lie in the Lorentz null cone.  
% Indeed, any~$e\in\Eigen(\pD_{\T^2};\sqrt{\zeta^2+m^2})$ satisfies
% \begin{align}
%     \bar{e}e=\frac{m}{\sqrt{m^2+\zeta^2}}|e|^2.
% \end{align}
% which is nonzero unless~$e=0$. 
% In particular,~$\vartheta_\infty$ cannot be in~$\Eigen(\pD_{\T^2};\sqrt{\zeta^2+m^2})$.  


Recall that~$\proj_{m,\zeta}$ is the~$L^2$-orthogonal projection to~$\Eigen(\pD_{\T^2};\sqrt{m^2+\zeta^2})$ and~$\proj_{m,\zeta}^\perp = \id -\proj_{m,\zeta}$ is the complementary projection.

% \begin{lemma}\label{lemma:null-cone-separation}
%     For any~$\nu\in (1,2)$, there exist~$\kappa_1>0$ and~$\kappa_2>0$ such that
%     \begin{align}\label{eq:kappa}
%         \|\vartheta\|_{L^4}=1,\quad \|\bar{\vartheta} \vartheta\|_{L^\nu}\leq\kappa_1 \quad\Longrightarrow\quad \|\proj_{m,\zeta}^\perp \vartheta\|_{L^4}\geq\kappa_2.
%     \end{align}
% \end{lemma}
% \begin{proof}
%     The corresponding statement for~$\nu\in(1,\frac{3}{2})$ was proved in~\cite{WZ2026Stationary}; but the argument works also for~$\nu\in (1,2)$.  
%     Since these constants are crucial for us, we give the detailed proof here. 
    
%     To see~\eqref{eq:kappa}, we argue by contradiction. 
%     Suppose  there exists a sequence~$(\vartheta_k)$ such that
%     \begin{align}
%         \|\vartheta_k\|_{L^4}=1, \quad \|\bar{\vartheta}_k\vartheta_k\|_{L^\nu} \to 0,
%         \quad \|\proj_{m,\zeta}^\perp \vartheta_k\|_{L^4}\to 0, \qquad \mbox{ as } \; k\to+\infty.
%     \end{align}
%     We observe that ~$\|\proj_{m,\zeta}\vartheta_k\|_{L^4}\leq C\|\vartheta_k\|_{L^4}=C$, with~$C$ depending on~$\T^2$ and the finite-dimensional space~$\Eigen(\pD_{\T^2};\sqrt{\zeta^2+m^2})$. 
%     Since any two norms on a finite-dimensional space are equivalent, the sequence~$(\proj_{m,\zeta}\vartheta_k)$ is bounded in~$\Eigen(\pD_{\T^2};\sqrt{m^2+\zeta^2})\cong\C^2$. 
%     Therefore, upon passing to a subsequence if necessary, the sequence~$\proj_{m,\zeta}\vartheta_k$ converges to some~$\vartheta_{*}$.
%     Together with the convergence~$\|\proj_{m,\zeta}^\perp \vartheta_k\|_{L^4}\to0$, we see that~$\vartheta_k$ converges in~$L^4$ to~$\vartheta_*$ with~$\|\vartheta_{*}\|_{L^4}=1$, and~$\bar{\vartheta}_k\vartheta_k$ converges in~$L^2$ to~$\bar{\vartheta}_*\vartheta_*$. 
%     Since~$\nu<2$, we see that~$\bar{\vartheta}_k\vartheta_k$ also converges in~$L^\nu$ to~$\bar{\vartheta}_*\vartheta_*$, hence~$\bar{\vartheta}_*\vartheta_*=0$ a.e. 
%     This is impossible, since~$\vartheta_*\in\Eigen(\pD_{\T^2};\sqrt{m^2+\zeta^2})$ and~$\bar{\vartheta}_*\vartheta_*= \frac{m}{\sqrt{\zeta^2 + m^2}}|\vartheta_*|^2$.
    
% \end{proof}
%       We remark that the above statement actually holds for any~$\nu\geq 1$ and there is no need for the requirement~$\nu<2$; but here we don't need that general form for the moment.

% \ 

%       We also need a resolvent estimate for~$\pD_{\T^2}-a $, for any~$a\in [0,\sqrt{\zeta^2+m^2}]$, on the orthogonal complement of the lowest positive eigenspace~$\Eigen(\pD_{\T^2};\sqrt{\zeta^2+m^2})$.

% \begin{lemma}\label{lemma:invertibility-orthogonal}
%     There exists a constant~$C=C(m,\zeta,\mathbb{T}^2)$ such that for any~$a\in [0,\sqrt{\zeta^2+m^2}]$, 
%     \begin{align}
%         \|\proj_{m,\zeta}^\perp \vartheta\|_{L^4}\leq C\|\proj_{m,\zeta}^\perp(\pD_{\T^2}-a)\vartheta\|_{L^{\frac43}}.
%     \end{align}
% \end{lemma}

% \begin{proof}
%     Suppose not, then there exist sequences~$(a_k)\subset [0,\sqrt{\zeta^2+m^2}]$ and~$(\vartheta_k)\subset W^{1,\frac{4}{3}}$ such that
%         \begin{align}
%             \|\proj_{m,\zeta}^\perp\vartheta_k\|_{L^4}=1, & &
%             \|\proj_{m,\zeta}^\perp(\pD_{\T^2}-a_k)\vartheta_k\|_{L^{\frac{4}{3}}}\to 0.
%         \end{align}
%         We may assume~$a_k\to a_*\in [0,\sqrt{\zeta^2+m^2}]$ by Heine--Borel theorem. 
%         Note that for any~$k$,
%         \begin{align}
%             \|\proj_{m,\zeta}^\perp \vartheta_k \|_{W^{1,\frac{4}{3}}}
%             \leq& C \parenthesis{ \| \proj_{m,\zeta}^\perp \vartheta_k\|_{L^{\frac{4}{3}}} + \|\pD_{\T^2}\proj_{m,\zeta}^\perp \vartheta_k\|_{L^{\frac{4}{3}}} } \\
%             \leq& C \parenthesis{
%             \| \proj_{m,\zeta}^\perp \vartheta_k\|_{L^{\frac{4}{3}}}
%             + \|(\pD_{\T^2}-a_k)\proj_{m,\zeta}^\perp \vartheta_k\|_{L^{\frac{4}{3}}} 
%             + a_k\|\proj_{m,\zeta}^\perp \vartheta_k\|_{L^{\frac{4}{3}}} } \\
%             \leq& C \parenthesis{ (\sqrt{m^2+\zeta^2}+1)\| \proj_{m,\zeta}^\perp \vartheta_k\|_{L^{\frac{4}{3}}} + \|\proj_{m,\zeta}^\perp (\pD_{\T^2}-a_k)\vartheta_k\|_{L^{\frac{4}{3}}} },
%         \end{align}
%         where~$C$ is independent of~$k$ and we have used the commutativity of~$\proj_{m,\zeta}^\perp$ with~$(\pD_{\T^2}-a_k)$ in the last inequality.
%         Hence the sequence~$(\proj_{m,\zeta}^\perp \vartheta_k)$ is also uniformly bounded in~$W^{1,\frac{4}{3}}$ and we may assume
%         \begin{align}
%             \proj_{m,\zeta}^\perp\vartheta_k \to \vartheta_*=\proj_{m,\zeta}^\perp \vartheta_*\quad  \mbox{ in } \;  L^{\frac{4}{3}}.
%         \end{align}
%         Moreover, from
%         \begin{align}
%             \pD_{\T^2}(\proj_{m,\zeta}^\perp\vartheta_k)= \proj_{m,\zeta}^\perp (\pD_{\T^2}-a_k)\vartheta_k+ a_k \proj_{m,\zeta}^\perp\vartheta_k \xrightarrow{L^{\frac{4}{3}}} 0+ a_*\vartheta_*
%         \end{align}
%         it follows that
%         \begin{align}
%             \pD_{\T^2}\vartheta_*=a_*\vartheta_*.
%         \end{align}
%         Since~$\vartheta_* =\proj_{m,\zeta}^\perp\vartheta_*\perp \Eigen(\pD_{\T^2};\sqrt{m^2+\zeta^2})$ while~$\dist(a_*, \Spect(\pD_{\T^2})\setminus\braces{\sqrt{m^2+\zeta^2}})>0$, we must have~$\vartheta_*=0$.
%         This in turn implies that
%         \begin{align}
%             \|\proj_{m,\zeta}^\perp \vartheta_k\|_{W^{1,\frac{4}{3}}} \to 0
%         \end{align}
%         as~$k\to+\infty$, contradicting~$\|\proj_{m,\zeta}^\perp\vartheta_k\|_{L^4}=1$ for any~$k$.
% \end{proof}
% The proof of Lemma~\ref{lemma:null-cone-separation} and Lemma~\ref{lemma:invertibility-orthogonal} will be deferred to next section. 

\ 

Now we have~$\|\bar{\vartheta}_n\vartheta_n\|_{L^\nu}\to 0$ and~$\|\vartheta_n\|_{L^4}=1$. 
By Lemma~\ref{lemma:null-cone-separation},
\begin{align}
    \|\proj_{m,\zeta}^\perp\vartheta_n\|_{L^4}\geq \kappa_2>0. 
\end{align}
On the other hand, using~\eqref{eq:normalized eqn}, (F5) and~\eqref{eq:uniform estimate for perturbation}, we have  
\begin{align}
    \| (\pD_{\T^2}-a_n)\vartheta_n\|_{L^{\frac{4}{3}}} 
    \leq& C(m,\zeta,\vol(\mathbb{T}^2),\beta,\alpha_2)\parenthesis{ \|\frac{1}{\lambda_n}\p F(\varphi_{\eps_n,a_n})\|_{L^{\frac{4\beta}{4+\beta}} } + \alpha_2\|\eps_n|\varphi_{\eps_n,a_n}|^{\alpha_2-2}\vartheta_n \|_{L^{\frac{\alpha_2}{\alpha_2-1}} } } \\
    \leq & C(m,\zeta,\vol(\mathbb{T}^2),\beta,\alpha_2)
    \parenthesis{ A_5\parenthesis{\delta+ C_\delta \parenthesis{ \int_{\mathbb{T}^2} F(\varphi_{\eps_n,a_n})\dv_{\T^2}}^{\frac{1}{\beta}} } + \frac{\eps_n^{\frac{1}{\alpha_2}}}{\lambda_n} } \\
    \leq &  C(m,\zeta,\vol(\mathbb{T}^2),\beta,\alpha_2, A_5)
    \parenthesis{ \delta + C_\delta (C^*(a_n))^{\frac{1}{\beta} } + \frac{\eps_n^{\frac{1}{\alpha_2}}}{\lambda_n} }.
\end{align}
Combining these two estimates using Lemma~\ref{lemma:invertibility-orthogonal} leads to 
\begin{align}
   0<\kappa_2\leq & \|\proj_{m,\zeta}^\perp \vartheta_n\|_{L^4}
    \leq C(m,\zeta) \|\proj_{m,\zeta}^\perp (\pD_{\T^2}-a_n)\vartheta_n\|_{L^{\frac{4}{3}}}  \\
    \leq & C(m,\zeta,\vol(\mathbb{T}^2),\beta,\alpha_2, A_5)
    \parenthesis{ \delta + C_\delta (C^*(a_n))^{\frac{1}{\beta} } + \frac{\eps_n^{\frac{1}{\alpha_2}}}{\lambda_n} }.
\end{align}
By first choosing~$\delta>0$ small, then sending~$n\to+\infty$ and recalling~$C^*(a_n)\to 0^+$ and~$\lambda_n\to+\infty$, we get a contradiction. 

This confirms the existence of~$m_*\in (\theta\sqrt{\zeta^2+m^2},\sqrt{\zeta^2+m^2})$ and~$C>0$ such that for any~$\eps\in (0,1]$ and any~$a\in (m_*,\sqrt{\zeta^2+m^2})$, we have 
\begin{align}
    \|\varphi_{\eps,a}\|_{L^4}\leq C. 
\end{align} 

\ 

\noindent\textbf{Step 3.} \emph{The perturbed solutions~$\varphi_{\eps,a}$ obtained as above are also uniformly bounded in~$H^1(\T^2;\C^4)$.}

With the above uniform~$L^4$ bound on~$\varphi_{\eps,a}$, we see that 
    \begin{align}
        |\pD_{\T^2}\varphi_{\eps,a}|
        =|a\varphi_{\eps,a} + \p F(\varphi_{\eps,a})+ \alpha_2\eps|\varphi_{\eps,a}|^{\alpha_2-2}\varphi_{\eps,a} |
        \leq
        C\parenthesis{ 1+ |\varphi_{\eps,a}|^{\alpha_2 -1}}
    \end{align}
    which is uniformly bounded in~$L^{\frac{4}{\alpha_2 -1 }}$ with~$\frac{4}{\alpha_2 -1}>\frac{4}{3}$.

    If~$\alpha_2\in (2,3]$, then~$\frac{4}{\alpha_2-1}\geq 2$, so the right-hand side of the Dirac equation is already uniformly bounded in~$L^2$. 
    Elliptic regularity immediately gives
    \begin{align}
        \|\varphi_{\eps,a}\|_{H^1}\leq C<+\infty, \quad \forall \eps\in (0,1].
    \end{align}
    If~$\alpha_2\in (3,4)$, then~$\frac{4}{\alpha_2-1}\in (\frac{4}{3},2)$, so the spinors~$\varphi_{\eps,a} $ are uniformly bounded in~$W^{1,\frac{4}{\alpha_2-1}} \hookrightarrow L^{\frac{4}{\alpha_2-3}}$ and~$\frac{4}{\alpha_2-3}>4$.
    We repeat this subcritical bootstrap procedure finitely many times until the exponent is sufficiently large so that the right-hand side belongs to~$L^2$. This ultimately yields a uniform~$H^1$-bound.

\end{proof}


\begin{thebibliography}{99}
\bibitem[]{Ammann2003Habil} Ammann, Bernd. A variational problem in conformal spin geometry. (2003)
\bibitem[]{BalabaneCazenave1988Existence} Balabane, Mikha\"el; Cazenave, Thierry; Douady, Adrien; Merle, Frank. Existence d'\'etats excit\'es pour une \'equation de {D}irac non lin\'eaire. C. R. Acad. Sci. Paris S\'er. I Math. 306 (1988) 117--120
\bibitem[]{BartschDing2006Solutions} Bartsch, Thomas; Ding, Yanheng. Solutions of nonlinear {D}irac equations. J. Differential Equations 226 (2006) 210--249
\bibitem[]{BartschXu2021spinorial} Bartsch, Thomas; Xu, Tian. A spinorial analogue of the {B}rezis-{N}irenberg theorem involving the critical {S}obolev exponent. J. Funct. Anal. 280 (2021) Paper No. 108991, 47
\bibitem[]{Cazenave1986Existence} Cazenave, Thierry; V\'azquez, Luis. Existence of localized solutions for a classical nonlinear {D}irac field. Comm. Math. Phys. 105 (1986) 35--47
\bibitem[]{DingLiXu2016bifurcation} Ding, Yanheng; Li, Jiongyue; Xu, Tian. Bifurcation on compact spin manifold. Calc. Var. Partial Differential Equations 55 (2016) Art. 90, 17
\bibitem[]{DingLiu2014periodic} Ding, Yanheng; Liu, Xiaoying. Periodic waves of nonlinear {D}irac equations. Nonlinear Anal. 109 (2014) 252--267
\bibitem[]{DingLiu2015Periodic} Ding, Yanheng; Liu, Xiaoying. Periodic solutions of a {D}irac equation with concave and convex nonlinearities. J. Differential Equations 258 (2015) 3567--3588
\bibitem[]{DingLiu2017Periodic} Ding, Yanheng; Liu, Xiaoying. Periodic solutions of an asymptotically linear {D}irac equation. Ann. Mat. Pura Appl. (4) 196 (2017) 717--735
\bibitem[]{DingRuf2008Solutions} Ding, Yanheng; Ruf, Bernhard. Solutions of a nonlinear {D}irac equation with external fields. Arch. Ration. Mech. Anal. 190 (2008) 57--82
\bibitem[]{DingXu2025localized} Ding, Yanheng; Xu, Tian. Localization approaches in strongly indefinite problems. 6 ([2025] \copyright 2025) xiii+281
\bibitem[]{Esteban2002overview} Esteban, Maria J.; S\'er\'e, Eric. An overview on linear and nonlinear {D}irac equations. Discrete Contin. Dyn. Syst. 8 (2002) 381--397
\bibitem[]{EstebanSere1995stationary} Esteban, Maria J.; S\'er\'e, \'Eric. Stationary states of the nonlinear {D}irac equation: a variational approach. Comm. Math. Phys. 171 (1995) 323--350
\bibitem[]{Finkelstein1951Non-linea} Finkelstein, R.; LeLevier, R.; Ruderman, M.. Non-linear spinor fields. Phys. Rev. (2) 83 (1951) 326--332
\bibitem[]{Friedrich1984zur} Friedrich, Thomas. Zur {A}bh\"angigkeit des {D}irac-operators von der {S}pin-{S}truktur. Colloq. Math. 48 (1984) 57--62
\bibitem[]{Ginoux2009Dirac} Ginoux, Nicolas. The {D}irac spectrum. 1976 (2009) xvi+156
\bibitem[]{Isobe2011nonlinear} Isobe, Takeshi. Nonlinear {D}irac equations with critical nonlinearities on compact {S}pin manifolds. J. Funct. Anal. 260 (2011) 253--307
\bibitem[]{Isobe2020Morse} Isobe, Takeshi. Morse homology for asymptotically linear {D}irac equations on compact manifolds. J. Differential Equations 269 (2020) 5062--5109
\bibitem[]{Lawson1989Spin} Lawson, Jr., H. Blaine; Michelsohn, Marie-Louise. Spin geometry. 38 (1989) xii+427
\bibitem[]{Merle1988Existence} Merle, F.. Existence of stationary states for nonlinear {D}irac equations. J. Differential Equations 74 (1988) 50--68
\bibitem[]{Ranada1983Classical} Ranada, A. F.. Classical nonlinear Dirac field models of extended particles. Math. Phys. Stud. 4 (1983) 271--291
\bibitem[]{SireXu2023variational} Sire, Yannick; Xu, Tian. A variational analysis of the spinorial {Y}amabe equation on product manifolds. Ann. Sc. Norm. Super. Pisa Cl. Sci. (5) 24 (2023) 205--248
\bibitem[]{Soler1970Classical} Soler, Mario. Classical, stable, nonlinear spinor field with positive rest energy. Phys. Rev. D 1 (1970) 2766--2769
\bibitem[]{Thaller1992Dirac} Thaller, Bernd. The {D}irac equation. (1992) xviii+357
\bibitem[]{Varilly2006Dirac} V{\'a}rilly, Joseph C.. Dirac Operators and Spectral Geometry. (2006)
\bibitem[]{WZ2026Stationary} Ruijun Wu; Fuping Zhang. Stationary periodic solutions to Nonlinear Dirac equations with non-coercive potentials. (2026)
\bibitem[]{YangDing2012Stationary} Yang, Minbo; Ding, Yanheng. Stationary states for nonlinear {D}irac equations with superlinear nonlinearities. Topol. Methods Nonlinear Anal. 39 (2012) 175--188
\end{thebibliography}
\end{document}
